
\documentclass{article}
\usepackage{colortbl}
\usepackage{makecell}
\usepackage{multirow}
\usepackage{supertabular}

\begin{document}

\newcounter{utterance}

\centering \large Interaction Transcript for game `clean\_up', experiment `1\_medium\_7obj\_ga', episode 1 with Qwen3.6{-}35B{-}A3B{-}FP8{-}without{-}reasoning.
\vspace{24pt}

{ \footnotesize  \setcounter{utterance}{1}
\setlength{\tabcolsep}{0pt}
\begin{supertabular}{c@{$\;$}|p{.15\linewidth}@{}p{.15\linewidth}p{.15\linewidth}p{.15\linewidth}p{.15\linewidth}p{.15\linewidth}}
    \# & $\;$A & \multicolumn{4}{c}{Game Master} & $\;\:$B\\
    \hline

    \theutterance \stepcounter{utterance}  
    & & \multicolumn{4}{p{0.6\linewidth}}{
        \cellcolor[rgb]{0.9,0.9,0.9}{
            \makecell[{{p{\linewidth}}}]{
                \texttt{\tiny{[P1$\langle$GM]}}
                \texttt{Is mise do mháistir cluiche, agus tá tú ag imirt cluiche comhoibríoch leis an ngreille seo a leanas mar chlár cluiche:} \\
\\ 
\texttt{```} \\
\texttt{ 1234567} \\
\texttt{╔══╤══╤═╗} \\
\texttt{║◌P│◌M│◌║ 1} \\
\texttt{╟──┼──┴─╢ 2} \\
\texttt{║◌◌│◌S◌O║ 3} \\
\texttt{╟──┘◌A◌◌║ 4} \\
\texttt{║D◌◌◌◌◌◌║ 5} \\
\texttt{║◌◌┌─┬──╢ 6} \\
\texttt{║T◌│◌│◌◌║ 7} \\
\texttt{╚══╧═╧══╝} \\
\\ 
\texttt{```} \\
\\ 
\texttt{* Taispeánann an imeall uachtarach comhordanáidí x ag méadú i dtreo na láimhe deise, agus an t{-}imeall ar dheis comhordanáidí y ag méadú síos.} \\
\texttt{* Cuirtear na rudaí seo a leanas go randamach ar do ghreille: 'P', 'O', 'T', 'S', 'D', 'A', 'M'.} \\
\\ 
\\ 
\texttt{Feiceann an t{-}imreoir eile leagan éagsúil den chlár cluiche, áit a gcuirtear na rudaí in áiteanna randamacha éagsúla. Ní féidir leat clár an imreora eile a fheiceáil, agus ní féidir leo do cheannsa a fheiceáil.} \\
\\ 
\\ 
\texttt{**Sprioc:**} \\
\\ 
\texttt{Caithfidh an dá imreoir na rudaí a bhogadh ar a gcúlra féin ionas go gcríochnóidh rudaí comhionanna ag na comhordanáidí céanna. Caithfidh tú cumarsáid a dhéanamh leis an imreoir eile chun cumraíocht sprice choiteann a aontú.} \\
\\ 
\\ 
\texttt{**Rialacha:**} \\
\\ 
\texttt{* I ngach babhta, is féidir leat ceann amháin go díreach den dá ordú seo a leanas a sheoladh:} \\
\texttt{ 1. `ABAIR: <TEACHTAIREACHT>`: chun teachtaireacht (gach rud suas go dtí an chéad bhriseadh líne eile) a sheoladh chuig an imreoir eile. Cuirfidh mé ar aghaidh í chuig do pháirtí.} \\
\texttt{2. `BOG: <RUD>, (<X>, <Y>)`: chun rud a bhogadh go suíomh nua, áit a bhfuil `<X>` mar an colún agus `<Y>` mar an ró. Cuirfidh mé in iúl duit an raibh do bhogadh rathúil nó nach raibh.} \\
\texttt{* Mura gcloíonn tú leis an bhformáid, nó má sheolann tú roinnt orduithe ag an am céanna, caithfidh mé pionós a ghearradh ort.} \\
\texttt{* Má charnaíonn an dá imreoir níos mó ná 16 pionóis, caillfidh sibh beirt an cluiche.} \\
\texttt{* Tá sé ríthábhachtach go ndéanann tú cumarsáid leis an imreoir eile maidir le do staid sprice! Is é an t{-}aon bhealach *amháin* ar féidir leat do straitéis a chur ar aghaidh chuig an imreoir eile ná an t{-}ordú `ABAIR: <TEACHTAIREACHT>` a úsáid!} \\
\\ 
\\ 
\texttt{**Rudaí a Bhogadh**} \\
\\ 
\texttt{* Ní féidir leat rudaí a bhogadh ach chuig cealla laistigh de theorainneacha na greille. Caithfidh an chill sprice a bheith folamh, is é sin, ní mór di ach an tsiombail '◌' a bheith inti.} \\
\texttt{* Má dhéanann tú iarracht rud a bhogadh go spota nach bhfuil folamh, nó má dhéanann tú iarracht é a bhogadh lasmuigh den ghreille, caithfidh mé pionós a ghearradh ort. Gheobhaidh tú triail eile.} \\
\texttt{* Sula ndéanann tú bogadh, seiceáil faoi dhó go bhfuil an spota sprice folamh, agus nach bhfuil aon litir, fráma ná líne ann!} \\
\\ 
\\ 
\texttt{**Deireadh an Chluiche**} \\
\\ 
\texttt{Má cheapann tú gur shroich tú an sprioc maidir le gach rud a ailíniú, is féidir leat iarraidh ar an imreoir eile an cluiche a chríochnú trí `ABAIR: críochnaithe?` a sheoladh. Má iarrann an t{-}imreoir eile ort an cluiche a chríochnú, agus má fhreagraíonn tú le `ABAIR: críochnaithe!`, tiocfaidh deireadh leis an gcluiche.} \\
\\ 
\texttt{Buaíonn an dá imreoir má chríochnaíonn an cluiche laistigh de 28 babhta, áit a shainmhínítear babhta amháin mar bheirt imreoirí ag seoladh ordú bailí an duine.} \\
\\ 
\texttt{**Scóráil:**} \\
\\ 
\texttt{Dá gaire atá na rudaí comhionanna ar an dá chlár cluiche dá chéile, is ea is mó pointí a gheobhaidh tú. Laghdaíonn pionóis do phointí. An féidir leat an taifead a shárú?} \\
\\ 
\\ 
\texttt{Seol teachtaireacht chuig an imreoir eile chun an cluiche a thosú!} \\
            }
        }
    }
    & \\ \\

    \theutterance \stepcounter{utterance}  
    & \multicolumn{4}{p{0.6\linewidth}}{
        \cellcolor[rgb]{0.9,0.9,0.9}{
            \makecell[{{p{\linewidth}}}]{
                \texttt{\tiny{[P1$\rangle$GM]}}
                \texttt{ABAIR: Dia duit! Seo a leanas mo staid reatha:} \\
\texttt{P ag (1,1)} \\
\texttt{M ag (2,1)} \\
\texttt{S ag (3,3)} \\
\texttt{O ag (4,3)} \\
\texttt{A ag (3,4)} \\
\texttt{D ag (1,5)} \\
\texttt{T ag (1,7)} \\
\\ 
\texttt{Cad iad do chomhordanáidí?} \\
            }
        }
    }
    & & \\ \\

    \theutterance \stepcounter{utterance}  
    & & & \multicolumn{2}{p{0.3\linewidth}}{
        \cellcolor[rgb]{0.9,0.9,0.9}{
            \makecell[{{p{\linewidth}}}]{
                \texttt{\tiny{[GM$|$GM]}}
                \texttt{ABAIR: Dia duit! Seo a leanas mo staid reatha:} \\
\texttt{P ag (1,1)} \\
\texttt{M ag (2,1)} \\
\texttt{S ag (3,3)} \\
\texttt{O ag (4,3)} \\
\texttt{A ag (3,4)} \\
\texttt{D ag (1,5)} \\
\texttt{T ag (1,7)} \\
\\ 
\texttt{Cad iad do chomhordanáidí?} \\
            }
        }
    }
    & & \\ \\

    \theutterance \stepcounter{utterance}  
    & & & \multicolumn{2}{p{0.3\linewidth}}{
        \cellcolor[rgb]{0.9,0.9,0.9}{
            \makecell[{{p{\linewidth}}}]{
                \texttt{\tiny{[GM$|$GM]}}
                \texttt{Invalid format: tail is not empty: '} \\
\texttt{P ag (1,1)} \\
\texttt{M ag (2,1)} \\
\texttt{S ag (3,3)} \\
\texttt{O ag (4,3)} \\
\texttt{A ag (3,4)} \\
\texttt{D ag (1,5)} \\
\texttt{T ag (1,7)} \\
\\ 
\texttt{Cad iad do chomhordanáidí?} \\
            }
        }
    }
    & & \\ \\

    \theutterance \stepcounter{utterance}  
    & & \multicolumn{4}{p{0.6\linewidth}}{
        \cellcolor[rgb]{0.9,0.9,0.9}{
            \makecell[{{p{\linewidth}}}]{
                \texttt{\tiny{[P1$\langle$GM]}}
                \texttt{Pionós: Ní cheadaítear aon rud a bheith i do theachtaireacht tar éis an ordaithe!} \\
\\ 
\texttt{Cinntigh nach bhfuil i do fhreagra ach `ABAIR: <TEACHTAIREACHT>` nó `BOG: <RUD>, (<X>, <Y>)`, agus tada eile!} \\
\\ 
\texttt{Tá 1 as 16 pionós carntha agaibh le chéile.} \\
\texttt{Déan iarracht arís, le do thoil!} \\
            }
        }
    }
    & \\ \\

    \theutterance \stepcounter{utterance}  
    & \multicolumn{4}{p{0.6\linewidth}}{
        \cellcolor[rgb]{0.9,0.9,0.9}{
            \makecell[{{p{\linewidth}}}]{
                \texttt{\tiny{[P1$\rangle$GM]}}
                \texttt{ABAIR: Dia duit! Seo a leanas mo staid reatha: P ag (1,1), M ag (2,1), S ag (3,3), O ag (4,3), A ag (3,4), D ag (1,5), T ag (1,7). Cad iad do chomhordanáidí?} \\
            }
        }
    }
    & & \\ \\

    \theutterance \stepcounter{utterance}  
    & & & \multicolumn{2}{p{0.3\linewidth}}{
        \cellcolor[rgb]{0.9,0.9,0.9}{
            \makecell[{{p{\linewidth}}}]{
                \texttt{\tiny{[GM$|$GM]}}
                \texttt{ABAIR: Dia duit! Seo a leanas mo staid reatha: P ag (1,1), M ag (2,1), S ag (3,3), O ag (4,3), A ag (3,4), D ag (1,5), T ag (1,7). Cad iad do chomhordanáidí?} \\
            }
        }
    }
    & & \\ \\

    \theutterance \stepcounter{utterance}  
    & & \multicolumn{4}{p{0.6\linewidth}}{
        \cellcolor[rgb]{0.9,0.9,0.9}{
            \makecell[{{p{\linewidth}}}]{
                \texttt{\tiny{[GM$\rangle$P2]}}
                \texttt{Is mise do mháistir cluiche, agus tá tú ag imirt cluiche comhoibríoch leis an ngreille seo a leanas mar chlár cluiche:} \\
\\ 
\texttt{```} \\
\texttt{ 1234567} \\
\texttt{╔══╤══╤═╗} \\
\texttt{║◌◌│◌◌│T║ 1} \\
\texttt{╟──┼──┴─╢ 2} \\
\texttt{║◌◌│◌S◌P║ 3} \\
\texttt{╟──┘◌◌MO║ 4} \\
\texttt{║D◌◌◌◌◌◌║ 5} \\
\texttt{║◌◌┌─┬──╢ 6} \\
\texttt{║◌◌│A│◌◌║ 7} \\
\texttt{╚══╧═╧══╝} \\
\\ 
\texttt{```} \\
\\ 
\texttt{* Taispeánann an imeall uachtarach comhordanáidí x ag méadú i dtreo na láimhe deise, agus an t{-}imeall ar dheis comhordanáidí y ag méadú síos.} \\
\texttt{* Cuirtear na rudaí seo a leanas go randamach ar do ghreille: 'P', 'O', 'T', 'S', 'D', 'A', 'M'.} \\
\\ 
\\ 
\texttt{Feiceann an t{-}imreoir eile leagan éagsúil den chlár cluiche, áit a gcuirtear na rudaí in áiteanna randamacha éagsúla. Ní féidir leat clár an imreora eile a fheiceáil, agus ní féidir leo do cheannsa a fheiceáil.} \\
\\ 
\\ 
\texttt{**Sprioc:**} \\
\\ 
\texttt{Caithfidh an dá imreoir na rudaí a bhogadh ar a gcúlra féin ionas go gcríochnóidh rudaí comhionanna ag na comhordanáidí céanna. Caithfidh tú cumarsáid a dhéanamh leis an imreoir eile chun cumraíocht sprice choiteann a aontú.} \\
\\ 
\\ 
\texttt{**Rialacha:**} \\
\\ 
\texttt{* I ngach babhta, is féidir leat ceann amháin go díreach den dá ordú seo a leanas a sheoladh:} \\
\texttt{ 1. `ABAIR: <TEACHTAIREACHT>`: chun teachtaireacht (gach rud suas go dtí an chéad bhriseadh líne eile) a sheoladh chuig an imreoir eile. Cuirfidh mé ar aghaidh í chuig do pháirtí.} \\
\texttt{2. `BOG: <RUD>, (<X>, <Y>)`: chun rud a bhogadh go suíomh nua, áit a bhfuil `<X>` mar an colún agus `<Y>` mar an ró. Cuirfidh mé in iúl duit an raibh do bhogadh rathúil nó nach raibh.} \\
\texttt{* Mura gcloíonn tú leis an bhformáid, nó má sheolann tú roinnt orduithe ag an am céanna, caithfidh mé pionós a ghearradh ort.} \\
\texttt{* Má charnaíonn an dá imreoir níos mó ná 16 pionóis, caillfidh sibh beirt an cluiche.} \\
\texttt{* Tá sé ríthábhachtach go ndéanann tú cumarsáid leis an imreoir eile maidir le do staid sprice! Is é an t{-}aon bhealach *amháin* ar féidir leat do straitéis a chur ar aghaidh chuig an imreoir eile ná an t{-}ordú `ABAIR: <TEACHTAIREACHT>` a úsáid!} \\
\\ 
\\ 
\texttt{**Rudaí a Bhogadh**} \\
\\ 
\texttt{* Ní féidir leat rudaí a bhogadh ach chuig cealla laistigh de theorainneacha na greille. Caithfidh an chill sprice a bheith folamh, is é sin, ní mór di ach an tsiombail '◌' a bheith inti.} \\
\texttt{* Má dhéanann tú iarracht rud a bhogadh go spota nach bhfuil folamh, nó má dhéanann tú iarracht é a bhogadh lasmuigh den ghreille, caithfidh mé pionós a ghearradh ort. Gheobhaidh tú triail eile.} \\
\texttt{* Sula ndéanann tú bogadh, seiceáil faoi dhó go bhfuil an spota sprice folamh, agus nach bhfuil aon litir, fráma ná líne ann!} \\
\\ 
\\ 
\texttt{**Deireadh an Chluiche**} \\
\\ 
\texttt{Má cheapann tú gur shroich tú an sprioc maidir le gach rud a ailíniú, is féidir leat iarraidh ar an imreoir eile an cluiche a chríochnú trí `ABAIR: críochnaithe?` a sheoladh. Má iarrann an t{-}imreoir eile ort an cluiche a chríochnú, agus má fhreagraíonn tú le `ABAIR: críochnaithe!`, tiocfaidh deireadh leis an gcluiche.} \\
\\ 
\texttt{Buaíonn an dá imreoir má chríochnaíonn an cluiche laistigh de 28 babhta, áit a shainmhínítear babhta amháin mar bheirt imreoirí ag seoladh ordú bailí an duine.} \\
\\ 
\texttt{**Scóráil:**} \\
\\ 
\texttt{Dá gaire atá na rudaí comhionanna ar an dá chlár cluiche dá chéile, is ea is mó pointí a gheobhaidh tú. Laghdaíonn pionóis do phointí. An féidir leat an taifead a shárú?} \\
\\ 
\texttt{Thosaigh an t{-}imreoir eile an cluiche trí an teachtaireacht seo a sheoladh:} \\
\\ 
\texttt{"Dia duit! Seo a leanas mo staid reatha: P ag (1,1), M ag (2,1), S ag (3,3), O ag (4,3), A ag (3,4), D ag (1,5), T ag (1,7). Cad iad do chomhordanáidí?"} \\
\\ 
\texttt{Cad é do chéad ordú?} \\
            }
        }
    }
    & \\ \\

    \theutterance \stepcounter{utterance}  
    & & & \multicolumn{4}{p{0.6\linewidth}}{
        \cellcolor[rgb]{0.9,0.9,0.9}{
            \makecell[{{p{\linewidth}}}]{
                \texttt{\tiny{[GM$\langle$P2]}}
                \texttt{ABAIR: Dia duit! Seo a leanas mo staid reatha: T ag (7,1), S ag (4,3), P ag (5,3), M ag (6,4), D ag (1,5), A ag (3,7).} \\
            }
        }
    }
     \\ \\

    \theutterance \stepcounter{utterance}  
    & & & \multicolumn{2}{p{0.3\linewidth}}{
        \cellcolor[rgb]{0.9,0.9,0.9}{
            \makecell[{{p{\linewidth}}}]{
                \texttt{\tiny{[GM$|$GM]}}
                \texttt{ABAIR: Dia duit! Seo a leanas mo staid reatha: T ag (7,1), S ag (4,3), P ag (5,3), M ag (6,4), D ag (1,5), A ag (3,7).} \\
            }
        }
    }
    & & \\ \\

    \theutterance \stepcounter{utterance}  
    & & \multicolumn{4}{p{0.6\linewidth}}{
        \cellcolor[rgb]{0.9,0.9,0.9}{
            \makecell[{{p{\linewidth}}}]{
                \texttt{\tiny{[P1$\langle$GM]}}
                \texttt{Cuireadh do theachtaireacht ar aghaidh chuig an imreoir eile.} \\
\\ 
\\ 
\texttt{Tá sibh ag imirt babhta 1 faoi láthair as uasmhéid 28.} \\
\texttt{Tá 1 as 16 pionós carntha agaibh le chéile.} \\
\texttt{Sheol an t{-}imreoir eile an teachtaireacht seo:} \\
\\ 
\\ 
\texttt{"Dia duit! Seo a leanas mo staid reatha: T ag (7,1), S ag (4,3), P ag (5,3), M ag (6,4), D ag (1,5), A ag (3,7)."} \\
\\ 
\\ 
\texttt{Cad é do chéad ordú eile?} \\
            }
        }
    }
    & \\ \\

    \theutterance \stepcounter{utterance}  
    & \multicolumn{4}{p{0.6\linewidth}}{
        \cellcolor[rgb]{0.9,0.9,0.9}{
            \makecell[{{p{\linewidth}}}]{
                \texttt{\tiny{[P1$\rangle$GM]}}
                \texttt{ABAIR: Go maith! Tá na comhordanáidí seo a leanas againn anois:} \\
\texttt{P: (1,1) vs (5,3)} \\
\texttt{M: (2,1) vs (6,4)} \\
\texttt{S: (3,3) vs (4,3)} \\
\texttt{O: (4,3) vs (gan)} \\
\texttt{A: (3,4) vs (3,7)} \\
\texttt{D: (1,5) vs (1,5) {-} Comhionann!} \\
\texttt{T: (1,7) vs (7,1)} \\
\\ 
\texttt{Tá D ag an suíomh cheart cheana féin.} \\
\texttt{Tá S ag (3,3) agam, agus ag an bpáirtí ag (4,3). Tá O agam ag (4,3) ach níl aon O aige. Mar sin, is féidir linn S agus O a mhalartú go héifeachtach má bhogaimid S go dtí (4,3) agus O go dtí (3,3). Ach ní mór dúinn a bheith cúramach.} \\
\\ 
\texttt{Straitéis:} \\
\texttt{1. Bog S ó (3,3) go (4,3) (áit a bhfuil S aige, ach tá O agam anseo, mar sin ní féidir liom S a bhogadh go dtí (4,3) go díreach mar tá O ann. Caithfidh mé O a bhogadh ar dtús).} \\
\texttt{2. Bog O ó (4,3) go (3,3) (áit a raibh S, ach tá S ag bogadh).} \\
\\ 
\texttt{Ní mór dúinn comhordú a dhéanamh.} \\
\texttt{Molaim:} \\
\texttt{1. Bog O ó (4,3) go (3,3) (folamh anois? Ní hea, tá S ann. Caithfidh mé S a bhogadh ar dtús nó O a bhogadh go háit eile).} \\
\\ 
\texttt{Ceartú:} \\
\texttt{Tá S ag (3,3) agam. Tá O ag (4,3) agam.} \\
\texttt{Tá S ag (4,3) aige.} \\
\texttt{Má bhogaim S ó (3,3) go háit eile, ansin bogaim O ó (4,3) go (3,3), ansin bogaim S go (4,3).} \\
\texttt{Cén áit inar féidir le S bogadh? (2,3) nó (3,2) nó (3,4) (tá A ann) nó (4,3) (tá O ann).} \\
\texttt{(2,3) atá folamh.} \\
\\ 
\texttt{Plan:} \\
\texttt{1. Bog S ó (3,3) go (2,3).} \\
\texttt{2. Bog O ó (4,3) go (3,3).} \\
\texttt{3. Bog S ó (2,3) go (4,3).} \\
\\ 
\texttt{Anois tá S ag (4,3) agam, agus tá S ag (4,3) aige. Comhionann!} \\
\texttt{Tá O ag (3,3) agam. Níl aon O aige. Caithfimid O a chur in áit éigin.} \\
\texttt{Tá P, M, A, T fágtha.} \\
\\ 
\texttt{P: (1,1) vs (5,3)} \\
\texttt{M: (2,1) vs (6,4)} \\
\texttt{A: (3,4) vs (3,7)} \\
\texttt{T: (1,7) vs (7,1)} \\
\\ 
\texttt{Tá D comhionann.} \\
\texttt{Tá S comhionann anois (tar éis na 3 bogadh).} \\
\\ 
\texttt{Anois, caithfimid P, M, A, T a ailíniú.} \\
\texttt{Tá P ag (1,1) agam. Tá P ag (5,3) aige.} \\
\texttt{Tá M ag (2,1) agam. Tá M ag (6,4) aige.} \\
\texttt{Tá A ag (3,4) agam. Tá A ag (3,7) aige.} \\
\texttt{Tá T ag (1,7) agam. Tá T ag (7,1) aige.} \\
\\ 
\texttt{Is fearr dúinn na rudaí a bhogadh go dtí na comhordanáidí céanna.} \\
\texttt{Molaim go n{-}úsáidfimid an chéad bhabhta chun S a shocrú mar a dúirt mé.} \\
\\ 
\texttt{BOG: S, (2,3)} \\
            }
        }
    }
    & & \\ \\

    \theutterance \stepcounter{utterance}  
    & & & \multicolumn{2}{p{0.3\linewidth}}{
        \cellcolor[rgb]{0.9,0.9,0.9}{
            \makecell[{{p{\linewidth}}}]{
                \texttt{\tiny{[GM$|$GM]}}
                \texttt{ABAIR: Go maith! Tá na comhordanáidí seo a leanas againn anois:} \\
\texttt{P: (1,1) vs (5,3)} \\
\texttt{M: (2,1) vs (6,4)} \\
\texttt{S: (3,3) vs (4,3)} \\
\texttt{O: (4,3) vs (gan)} \\
\texttt{A: (3,4) vs (3,7)} \\
\texttt{D: (1,5) vs (1,5) {-} Comhionann!} \\
\texttt{T: (1,7) vs (7,1)} \\
\\ 
\texttt{Tá D ag an suíomh cheart cheana féin.} \\
\texttt{Tá S ag (3,3) agam, agus ag an bpáirtí ag (4,3). Tá O agam ag (4,3) ach níl aon O aige. Mar sin, is féidir linn S agus O a mhalartú go héifeachtach má bhogaimid S go dtí (4,3) agus O go dtí (3,3). Ach ní mór dúinn a bheith cúramach.} \\
\\ 
\texttt{Straitéis:} \\
\texttt{1. Bog S ó (3,3) go (4,3) (áit a bhfuil S aige, ach tá O agam anseo, mar sin ní féidir liom S a bhogadh go dtí (4,3) go díreach mar tá O ann. Caithfidh mé O a bhogadh ar dtús).} \\
\texttt{2. Bog O ó (4,3) go (3,3) (áit a raibh S, ach tá S ag bogadh).} \\
\\ 
\texttt{Ní mór dúinn comhordú a dhéanamh.} \\
\texttt{Molaim:} \\
\texttt{1. Bog O ó (4,3) go (3,3) (folamh anois? Ní hea, tá S ann. Caithfidh mé S a bhogadh ar dtús nó O a bhogadh go háit eile).} \\
\\ 
\texttt{Ceartú:} \\
\texttt{Tá S ag (3,3) agam. Tá O ag (4,3) agam.} \\
\texttt{Tá S ag (4,3) aige.} \\
\texttt{Má bhogaim S ó (3,3) go háit eile, ansin bogaim O ó (4,3) go (3,3), ansin bogaim S go (4,3).} \\
\texttt{Cén áit inar féidir le S bogadh? (2,3) nó (3,2) nó (3,4) (tá A ann) nó (4,3) (tá O ann).} \\
\texttt{(2,3) atá folamh.} \\
\\ 
\texttt{Plan:} \\
\texttt{1. Bog S ó (3,3) go (2,3).} \\
\texttt{2. Bog O ó (4,3) go (3,3).} \\
\texttt{3. Bog S ó (2,3) go (4,3).} \\
\\ 
\texttt{Anois tá S ag (4,3) agam, agus tá S ag (4,3) aige. Comhionann!} \\
\texttt{Tá O ag (3,3) agam. Níl aon O aige. Caithfimid O a chur in áit éigin.} \\
\texttt{Tá P, M, A, T fágtha.} \\
\\ 
\texttt{P: (1,1) vs (5,3)} \\
\texttt{M: (2,1) vs (6,4)} \\
\texttt{A: (3,4) vs (3,7)} \\
\texttt{T: (1,7) vs (7,1)} \\
\\ 
\texttt{Tá D comhionann.} \\
\texttt{Tá S comhionann anois (tar éis na 3 bogadh).} \\
\\ 
\texttt{Anois, caithfimid P, M, A, T a ailíniú.} \\
\texttt{Tá P ag (1,1) agam. Tá P ag (5,3) aige.} \\
\texttt{Tá M ag (2,1) agam. Tá M ag (6,4) aige.} \\
\texttt{Tá A ag (3,4) agam. Tá A ag (3,7) aige.} \\
\texttt{Tá T ag (1,7) agam. Tá T ag (7,1) aige.} \\
\\ 
\texttt{Is fearr dúinn na rudaí a bhogadh go dtí na comhordanáidí céanna.} \\
\texttt{Molaim go n{-}úsáidfimid an chéad bhabhta chun S a shocrú mar a dúirt mé.} \\
\\ 
\texttt{BOG: S, (2,3)} \\
            }
        }
    }
    & & \\ \\

    \theutterance \stepcounter{utterance}  
    & & & \multicolumn{2}{p{0.3\linewidth}}{
        \cellcolor[rgb]{0.9,0.9,0.9}{
            \makecell[{{p{\linewidth}}}]{
                \texttt{\tiny{[GM$|$GM]}}
                \texttt{Invalid response format: ABAIR: Go maith! Tá na comhordanáidí seo a leanas againn anois:} \\
\texttt{P: (1,1) vs (5,3)} \\
\texttt{M: (2,1) vs (6,4)} \\
\texttt{S: (3,3) vs (4,3)} \\
\texttt{O: (4,3) vs (gan)} \\
\texttt{A: (3,4) vs (3,7)} \\
\texttt{D: (1,5) vs (1,5) {-} Comhionann!} \\
\texttt{T: (1,7) vs (7,1)} \\
\\ 
\texttt{Tá D ag an suíomh cheart cheana féin.} \\
\texttt{Tá S ag (3,3) agam, agus ag an bpáirtí ag (4,3). Tá O agam ag (4,3) ach níl aon O aige. Mar sin, is féidir linn S agus O a mhalartú go héifeachtach má bhogaimid S go dtí (4,3) agus O go dtí (3,3). Ach ní mór dúinn a bheith cúramach.} \\
\\ 
\texttt{Straitéis:} \\
\texttt{1. Bog S ó (3,3) go (4,3) (áit a bhfuil S aige, ach tá O agam anseo, mar sin ní féidir liom S a bhogadh go dtí (4,3) go díreach mar tá O ann. Caithfidh mé O a bhogadh ar dtús).} \\
\texttt{2. Bog O ó (4,3) go (3,3) (áit a raibh S, ach tá S ag bogadh).} \\
\\ 
\texttt{Ní mór dúinn comhordú a dhéanamh.} \\
\texttt{Molaim:} \\
\texttt{1. Bog O ó (4,3) go (3,3) (folamh anois? Ní hea, tá S ann. Caithfidh mé S a bhogadh ar dtús nó O a bhogadh go háit eile).} \\
\\ 
\texttt{Ceartú:} \\
\texttt{Tá S ag (3,3) agam. Tá O ag (4,3) agam.} \\
\texttt{Tá S ag (4,3) aige.} \\
\texttt{Má bhogaim S ó (3,3) go háit eile, ansin bogaim O ó (4,3) go (3,3), ansin bogaim S go (4,3).} \\
\texttt{Cén áit inar féidir le S bogadh? (2,3) nó (3,2) nó (3,4) (tá A ann) nó (4,3) (tá O ann).} \\
\texttt{(2,3) atá folamh.} \\
\\ 
\texttt{Plan:} \\
\texttt{1. Bog S ó (3,3) go (2,3).} \\
\texttt{2. Bog O ó (4,3) go (3,3).} \\
\texttt{3. Bog S ó (2,3) go (4,3).} \\
\\ 
\texttt{Anois tá S ag (4,3) agam, agus tá S ag (4,3) aige. Comhionann!} \\
\texttt{Tá O ag (3,3) agam. Níl aon O aige. Caithfimid O a chur in áit éigin.} \\
\texttt{Tá P, M, A, T fágtha.} \\
\\ 
\texttt{P: (1,1) vs (5,3)} \\
\texttt{M: (2,1) vs (6,4)} \\
\texttt{A: (3,4) vs (3,7)} \\
\texttt{T: (1,7) vs (7,1)} \\
\\ 
\texttt{Tá D comhionann.} \\
\texttt{Tá S comhionann anois (tar éis na 3 bogadh).} \\
\\ 
\texttt{Anois, caithfimid P, M, A, T a ailíniú.} \\
\texttt{Tá P ag (1,1) agam. Tá P ag (5,3) aige.} \\
\texttt{Tá M ag (2,1) agam. Tá M ag (6,4) aige.} \\
\texttt{Tá A ag (3,4) agam. Tá A ag (3,7) aige.} \\
\texttt{Tá T ag (1,7) agam. Tá T ag (7,1) aige.} \\
\\ 
\texttt{Is fearr dúinn na rudaí a bhogadh go dtí na comhordanáidí céanna.} \\
\texttt{Molaim go n{-}úsáidfimid an chéad bhabhta chun S a shocrú mar a dúirt mé.} \\
\\ 
\texttt{BOG: S, (2,3)} \\
            }
        }
    }
    & & \\ \\

    \theutterance \stepcounter{utterance}  
    & & \multicolumn{4}{p{0.6\linewidth}}{
        \cellcolor[rgb]{0.9,0.9,0.9}{
            \makecell[{{p{\linewidth}}}]{
                \texttt{\tiny{[P1$\langle$GM]}}
                \texttt{Pionós: Tá níos mó ná ordú amháin i do theachtaireacht!} \\
\\ 
\texttt{Cinntigh nach bhfuil i do fhreagra ach `ABAIR: <TEACHTAIREACHT>` nó `BOG: <RUD>, (<X>, <Y>)`, agus tada eile!} \\
\\ 
\texttt{Tá 2 as 16 pionós carntha agaibh le chéile.} \\
\texttt{Déan iarracht arís, le do thoil!} \\
            }
        }
    }
    & \\ \\

    \theutterance \stepcounter{utterance}  
    & \multicolumn{4}{p{0.6\linewidth}}{
        \cellcolor[rgb]{0.9,0.9,0.9}{
            \makecell[{{p{\linewidth}}}]{
                \texttt{\tiny{[P1$\rangle$GM]}}
                \texttt{ABAIR: Tá an plean agam S a bhogadh go (2,3) ar dtús chun an bealach a ghlanadh. An bhfuil tú ar aon intinn leis?} \\
            }
        }
    }
    & & \\ \\

    \theutterance \stepcounter{utterance}  
    & & & \multicolumn{2}{p{0.3\linewidth}}{
        \cellcolor[rgb]{0.9,0.9,0.9}{
            \makecell[{{p{\linewidth}}}]{
                \texttt{\tiny{[GM$|$GM]}}
                \texttt{ABAIR: Tá an plean agam S a bhogadh go (2,3) ar dtús chun an bealach a ghlanadh. An bhfuil tú ar aon intinn leis?} \\
            }
        }
    }
    & & \\ \\

    \theutterance \stepcounter{utterance}  
    & & \multicolumn{4}{p{0.6\linewidth}}{
        \cellcolor[rgb]{0.9,0.9,0.9}{
            \makecell[{{p{\linewidth}}}]{
                \texttt{\tiny{[GM$\rangle$P2]}}
                \texttt{Cuireadh do theachtaireacht ar aghaidh chuig an imreoir eile.} \\
\\ 
\\ 
\texttt{Tá sibh ag imirt babhta 2 faoi láthair as uasmhéid 28.} \\
\texttt{Tá 2 as 16 pionós carntha agaibh le chéile.} \\
\texttt{Sheol an t{-}imreoir eile an teachtaireacht seo:} \\
\\ 
\\ 
\texttt{"Tá an plean agam S a bhogadh go (2,3) ar dtús chun an bealach a ghlanadh. An bhfuil tú ar aon intinn leis?"} \\
\\ 
\\ 
\texttt{Cad é do chéad ordú eile?} \\
            }
        }
    }
    & \\ \\

    \theutterance \stepcounter{utterance}  
    & & & \multicolumn{4}{p{0.6\linewidth}}{
        \cellcolor[rgb]{0.9,0.9,0.9}{
            \makecell[{{p{\linewidth}}}]{
                \texttt{\tiny{[GM$\langle$P2]}}
                \texttt{ABAIR: Táim sásta leis sin. Bog S go (2,3). Ansin, ba chóir dúinn P a bhogadh go (2,1) chun é a ailíniú le do chuid.} \\
            }
        }
    }
     \\ \\

    \theutterance \stepcounter{utterance}  
    & & & \multicolumn{2}{p{0.3\linewidth}}{
        \cellcolor[rgb]{0.9,0.9,0.9}{
            \makecell[{{p{\linewidth}}}]{
                \texttt{\tiny{[GM$|$GM]}}
                \texttt{ABAIR: Táim sásta leis sin. Bog S go (2,3). Ansin, ba chóir dúinn P a bhogadh go (2,1) chun é a ailíniú le do chuid.} \\
            }
        }
    }
    & & \\ \\

    \theutterance \stepcounter{utterance}  
    & & \multicolumn{4}{p{0.6\linewidth}}{
        \cellcolor[rgb]{0.9,0.9,0.9}{
            \makecell[{{p{\linewidth}}}]{
                \texttt{\tiny{[P1$\langle$GM]}}
                \texttt{Cuireadh do theachtaireacht ar aghaidh chuig an imreoir eile.} \\
\\ 
\\ 
\texttt{Tá sibh ag imirt babhta 2 faoi láthair as uasmhéid 28.} \\
\texttt{Tá 2 as 16 pionós carntha agaibh le chéile.} \\
\texttt{Sheol an t{-}imreoir eile an teachtaireacht seo:} \\
\\ 
\\ 
\texttt{"Táim sásta leis sin. Bog S go (2,3). Ansin, ba chóir dúinn P a bhogadh go (2,1) chun é a ailíniú le do chuid."} \\
\\ 
\\ 
\texttt{Cad é do chéad ordú eile?} \\
            }
        }
    }
    & \\ \\

    \theutterance \stepcounter{utterance}  
    & \multicolumn{4}{p{0.6\linewidth}}{
        \cellcolor[rgb]{0.9,0.9,0.9}{
            \makecell[{{p{\linewidth}}}]{
                \texttt{\tiny{[P1$\rangle$GM]}}
                \texttt{BOG: S, (2,3)} \\
            }
        }
    }
    & & \\ \\

    \theutterance \stepcounter{utterance}  
    & & & \multicolumn{2}{p{0.3\linewidth}}{
        \cellcolor[rgb]{0.9,0.9,0.9}{
            \makecell[{{p{\linewidth}}}]{
                \texttt{\tiny{[GM$|$GM]}}
                \texttt{BOG: S, (2,3)} \\
            }
        }
    }
    & & \\ \\

    \theutterance \stepcounter{utterance}  
    & & & \multicolumn{2}{p{0.3\linewidth}}{
        \cellcolor[rgb]{0.9,0.9,0.9}{
            \makecell[{{p{\linewidth}}}]{
                \texttt{\tiny{[GM$|$GM]}}
                \texttt{S moved to (2, 3)} \\
            }
        }
    }
    & & \\ \\

    \theutterance \stepcounter{utterance}  
    & & \multicolumn{4}{p{0.6\linewidth}}{
        \cellcolor[rgb]{0.9,0.9,0.9}{
            \makecell[{{p{\linewidth}}}]{
                \texttt{\tiny{[GM$\rangle$P2]}}
                \texttt{Cuireadh do theachtaireacht ar aghaidh chuig an imreoir eile.} \\
\\ 
\\ 
\texttt{Tá sibh ag imirt babhta 3 faoi láthair as uasmhéid 28.} \\
\texttt{Tá 2 as 16 pionós carntha agaibh le chéile.} \\
\texttt{Bhog an t{-}imreoir eile réad ar a ngreille.} \\
\\ 
\\ 
\texttt{Cad é do chéad ordú eile?} \\
            }
        }
    }
    & \\ \\

    \theutterance \stepcounter{utterance}  
    & & & \multicolumn{4}{p{0.6\linewidth}}{
        \cellcolor[rgb]{0.9,0.9,0.9}{
            \makecell[{{p{\linewidth}}}]{
                \texttt{\tiny{[GM$\langle$P2]}}
                \texttt{BOG: S, (2,3)} \\
            }
        }
    }
     \\ \\

    \theutterance \stepcounter{utterance}  
    & & & \multicolumn{2}{p{0.3\linewidth}}{
        \cellcolor[rgb]{0.9,0.9,0.9}{
            \makecell[{{p{\linewidth}}}]{
                \texttt{\tiny{[GM$|$GM]}}
                \texttt{BOG: S, (2,3)} \\
            }
        }
    }
    & & \\ \\

    \theutterance \stepcounter{utterance}  
    & & & \multicolumn{2}{p{0.3\linewidth}}{
        \cellcolor[rgb]{0.9,0.9,0.9}{
            \makecell[{{p{\linewidth}}}]{
                \texttt{\tiny{[GM$|$GM]}}
                \texttt{S moved to (2, 3)} \\
            }
        }
    }
    & & \\ \\

    \theutterance \stepcounter{utterance}  
    & & \multicolumn{4}{p{0.6\linewidth}}{
        \cellcolor[rgb]{0.9,0.9,0.9}{
            \makecell[{{p{\linewidth}}}]{
                \texttt{\tiny{[P1$\langle$GM]}}
                \texttt{Bogadh 'S' go (2,3) go rathúil. Seo mar atá do ghreille nuashonraithe:} \\
\texttt{```} \\
\texttt{ 1234567} \\
\texttt{╔══╤══╤═╗} \\
\texttt{║◌P│◌M│◌║ 1} \\
\texttt{╟──┼──┴─╢ 2} \\
\texttt{║◌S│◌◌◌O║ 3} \\
\texttt{╟──┘◌A◌◌║ 4} \\
\texttt{║D◌◌◌◌◌◌║ 5} \\
\texttt{║◌◌┌─┬──╢ 6} \\
\texttt{║T◌│◌│◌◌║ 7} \\
\texttt{╚══╧═╧══╝} \\
\texttt{```} \\
\texttt{Tá sibh ag imirt babhta 3 faoi láthair as uasmhéid 28.} \\
\texttt{Tá 2 as 16 pionós carntha agaibh le chéile.} \\
\texttt{Bhog an t{-}imreoir eile réad ar a ngreille.} \\
\\ 
\\ 
\texttt{Cad é do chéad ordú eile?} \\
            }
        }
    }
    & \\ \\

    \theutterance \stepcounter{utterance}  
    & \multicolumn{4}{p{0.6\linewidth}}{
        \cellcolor[rgb]{0.9,0.9,0.9}{
            \makecell[{{p{\linewidth}}}]{
                \texttt{\tiny{[P1$\rangle$GM]}}
                \texttt{ABAIR: Go raibh maith agat. Anois, de réir do mholta, ba chóir dúinn P a bhogadh go (2,1). Ach tá M ag (2,1) agam. Ní féidir liom P a bhogadh go dtí áit atá ocupáilte. An bhfuil P agat ag (5,3)? Má tá, caithfimid P a bhogadh go dtí suíomh atá folamh ar mo chlár chun é a ailíniú leat. Molaim go n{-}úsáidfimid (1,2) nó (2,2) mar sprioc do P, ach caithfimid a bheith cinnte go bhfuil sé folamh. Ar dtús, bogaimid M ó (2,1) go dtí (1,1) (áit a raibh P, ach tá P ann fós). Ní féidir.} \\
\\ 
\texttt{Ceartú: Tá P ag (1,1) agus M ag (2,1).} \\
\texttt{Má bhogaim P go (1,2) (folamh), ansin is féidir le M bogadh go (1,1).} \\
\texttt{Ach deir an páirtí "Bog P go (2,1)". Ciallaíonn sé sin go gcaithfidh M bogadh ar dtús.} \\
\texttt{Molaim:} \\
\texttt{1. Bog M ó (2,1) go (1,1) {-} NÍ FÉIDHIR, tá P ann.} \\
\texttt{2. Bog P ó (1,1) go (1,2) {-} Folamh.} \\
\texttt{3. Ansin bog M ó (2,1) go (1,1).} \\
\texttt{4. Ansin bog P ó (1,2) go (2,1).} \\
\\ 
\texttt{Ach tá sé seo casta. Ina ionad sin, má tá P agat ag (5,3), agus tá P agam ag (1,1), is fearr dúinn comhordanáidí aontaithe a úsáid.} \\
\texttt{Molaim go n{-}úsáidfimid (2,1) mar sprioc do P.} \\
\texttt{Chun é sin a dhéanamh, caithfidh M bogadh.} \\
\texttt{Bog M ó (2,1) go (2,2) (folamh).} \\
\texttt{Ansin bog P ó (1,1) go (2,1).} \\
\\ 
\texttt{An bhfuil tú sásta leis an bplean seo?} \\
\texttt{1. Bog M go (2,2).} \\
\texttt{2. Bog P go (2,1).} \\
\\ 
\texttt{Seolfaidh mé an chéad bhogadh anois.} \\
\\ 
\texttt{BOG: M, (2,2)} \\
            }
        }
    }
    & & \\ \\

    \theutterance \stepcounter{utterance}  
    & & & \multicolumn{2}{p{0.3\linewidth}}{
        \cellcolor[rgb]{0.9,0.9,0.9}{
            \makecell[{{p{\linewidth}}}]{
                \texttt{\tiny{[GM$|$GM]}}
                \texttt{ABAIR: Go raibh maith agat. Anois, de réir do mholta, ba chóir dúinn P a bhogadh go (2,1). Ach tá M ag (2,1) agam. Ní féidir liom P a bhogadh go dtí áit atá ocupáilte. An bhfuil P agat ag (5,3)? Má tá, caithfimid P a bhogadh go dtí suíomh atá folamh ar mo chlár chun é a ailíniú leat. Molaim go n{-}úsáidfimid (1,2) nó (2,2) mar sprioc do P, ach caithfimid a bheith cinnte go bhfuil sé folamh. Ar dtús, bogaimid M ó (2,1) go dtí (1,1) (áit a raibh P, ach tá P ann fós). Ní féidir.} \\
\\ 
\texttt{Ceartú: Tá P ag (1,1) agus M ag (2,1).} \\
\texttt{Má bhogaim P go (1,2) (folamh), ansin is féidir le M bogadh go (1,1).} \\
\texttt{Ach deir an páirtí "Bog P go (2,1)". Ciallaíonn sé sin go gcaithfidh M bogadh ar dtús.} \\
\texttt{Molaim:} \\
\texttt{1. Bog M ó (2,1) go (1,1) {-} NÍ FÉIDHIR, tá P ann.} \\
\texttt{2. Bog P ó (1,1) go (1,2) {-} Folamh.} \\
\texttt{3. Ansin bog M ó (2,1) go (1,1).} \\
\texttt{4. Ansin bog P ó (1,2) go (2,1).} \\
\\ 
\texttt{Ach tá sé seo casta. Ina ionad sin, má tá P agat ag (5,3), agus tá P agam ag (1,1), is fearr dúinn comhordanáidí aontaithe a úsáid.} \\
\texttt{Molaim go n{-}úsáidfimid (2,1) mar sprioc do P.} \\
\texttt{Chun é sin a dhéanamh, caithfidh M bogadh.} \\
\texttt{Bog M ó (2,1) go (2,2) (folamh).} \\
\texttt{Ansin bog P ó (1,1) go (2,1).} \\
\\ 
\texttt{An bhfuil tú sásta leis an bplean seo?} \\
\texttt{1. Bog M go (2,2).} \\
\texttt{2. Bog P go (2,1).} \\
\\ 
\texttt{Seolfaidh mé an chéad bhogadh anois.} \\
\\ 
\texttt{BOG: M, (2,2)} \\
            }
        }
    }
    & & \\ \\

    \theutterance \stepcounter{utterance}  
    & & & \multicolumn{2}{p{0.3\linewidth}}{
        \cellcolor[rgb]{0.9,0.9,0.9}{
            \makecell[{{p{\linewidth}}}]{
                \texttt{\tiny{[GM$|$GM]}}
                \texttt{Invalid response format: ABAIR: Go raibh maith agat. Anois, de réir do mholta, ba chóir dúinn P a bhogadh go (2,1). Ach tá M ag (2,1) agam. Ní féidir liom P a bhogadh go dtí áit atá ocupáilte. An bhfuil P agat ag (5,3)? Má tá, caithfimid P a bhogadh go dtí suíomh atá folamh ar mo chlár chun é a ailíniú leat. Molaim go n{-}úsáidfimid (1,2) nó (2,2) mar sprioc do P, ach caithfimid a bheith cinnte go bhfuil sé folamh. Ar dtús, bogaimid M ó (2,1) go dtí (1,1) (áit a raibh P, ach tá P ann fós). Ní féidir.} \\
\\ 
\texttt{Ceartú: Tá P ag (1,1) agus M ag (2,1).} \\
\texttt{Má bhogaim P go (1,2) (folamh), ansin is féidir le M bogadh go (1,1).} \\
\texttt{Ach deir an páirtí "Bog P go (2,1)". Ciallaíonn sé sin go gcaithfidh M bogadh ar dtús.} \\
\texttt{Molaim:} \\
\texttt{1. Bog M ó (2,1) go (1,1) {-} NÍ FÉIDHIR, tá P ann.} \\
\texttt{2. Bog P ó (1,1) go (1,2) {-} Folamh.} \\
\texttt{3. Ansin bog M ó (2,1) go (1,1).} \\
\texttt{4. Ansin bog P ó (1,2) go (2,1).} \\
\\ 
\texttt{Ach tá sé seo casta. Ina ionad sin, má tá P agat ag (5,3), agus tá P agam ag (1,1), is fearr dúinn comhordanáidí aontaithe a úsáid.} \\
\texttt{Molaim go n{-}úsáidfimid (2,1) mar sprioc do P.} \\
\texttt{Chun é sin a dhéanamh, caithfidh M bogadh.} \\
\texttt{Bog M ó (2,1) go (2,2) (folamh).} \\
\texttt{Ansin bog P ó (1,1) go (2,1).} \\
\\ 
\texttt{An bhfuil tú sásta leis an bplean seo?} \\
\texttt{1. Bog M go (2,2).} \\
\texttt{2. Bog P go (2,1).} \\
\\ 
\texttt{Seolfaidh mé an chéad bhogadh anois.} \\
\\ 
\texttt{BOG: M, (2,2)} \\
            }
        }
    }
    & & \\ \\

    \theutterance \stepcounter{utterance}  
    & & \multicolumn{4}{p{0.6\linewidth}}{
        \cellcolor[rgb]{0.9,0.9,0.9}{
            \makecell[{{p{\linewidth}}}]{
                \texttt{\tiny{[P1$\langle$GM]}}
                \texttt{Pionós: Tá níos mó ná ordú amháin i do theachtaireacht!} \\
\\ 
\texttt{Cinntigh nach bhfuil i do fhreagra ach `ABAIR: <TEACHTAIREACHT>` nó `BOG: <RUD>, (<X>, <Y>)`, agus tada eile!} \\
\\ 
\texttt{Tá 3 as 16 pionós carntha agaibh le chéile.} \\
\texttt{Déan iarracht arís, le do thoil!} \\
            }
        }
    }
    & \\ \\

    \theutterance \stepcounter{utterance}  
    & \multicolumn{4}{p{0.6\linewidth}}{
        \cellcolor[rgb]{0.9,0.9,0.9}{
            \makecell[{{p{\linewidth}}}]{
                \texttt{\tiny{[P1$\rangle$GM]}}
                \texttt{BOG: M, (2,2)} \\
            }
        }
    }
    & & \\ \\

    \theutterance \stepcounter{utterance}  
    & & & \multicolumn{2}{p{0.3\linewidth}}{
        \cellcolor[rgb]{0.9,0.9,0.9}{
            \makecell[{{p{\linewidth}}}]{
                \texttt{\tiny{[GM$|$GM]}}
                \texttt{BOG: M, (2,2)} \\
            }
        }
    }
    & & \\ \\

    \theutterance \stepcounter{utterance}  
    & & & \multicolumn{2}{p{0.3\linewidth}}{
        \cellcolor[rgb]{0.9,0.9,0.9}{
            \makecell[{{p{\linewidth}}}]{
                \texttt{\tiny{[GM$|$GM]}}
                \texttt{Pionós: níl (2,2) folamh, ach tá '─' ann.} \\
\\ 
\texttt{Tá 4 as 16 pionós carntha agaibh le chéile.} \\
\texttt{Déan iarracht arís, le do thoil!} \\
            }
        }
    }
    & & \\ \\

    \theutterance \stepcounter{utterance}  
    & & \multicolumn{4}{p{0.6\linewidth}}{
        \cellcolor[rgb]{0.9,0.9,0.9}{
            \makecell[{{p{\linewidth}}}]{
                \texttt{\tiny{[P1$\langle$GM]}}
                \texttt{Pionós: níl (2,2) folamh, ach tá '─' ann.} \\
\\ 
\texttt{Tá 4 as 16 pionós carntha agaibh le chéile.} \\
\texttt{Déan iarracht arís, le do thoil!} \\
            }
        }
    }
    & \\ \\

    \theutterance \stepcounter{utterance}  
    & \multicolumn{4}{p{0.6\linewidth}}{
        \cellcolor[rgb]{0.9,0.9,0.9}{
            \makecell[{{p{\linewidth}}}]{
                \texttt{\tiny{[P1$\rangle$GM]}}
                \texttt{BOG: M, (1,2)} \\
            }
        }
    }
    & & \\ \\

    \theutterance \stepcounter{utterance}  
    & & & \multicolumn{2}{p{0.3\linewidth}}{
        \cellcolor[rgb]{0.9,0.9,0.9}{
            \makecell[{{p{\linewidth}}}]{
                \texttt{\tiny{[GM$|$GM]}}
                \texttt{BOG: M, (1,2)} \\
            }
        }
    }
    & & \\ \\

    \theutterance \stepcounter{utterance}  
    & & & \multicolumn{2}{p{0.3\linewidth}}{
        \cellcolor[rgb]{0.9,0.9,0.9}{
            \makecell[{{p{\linewidth}}}]{
                \texttt{\tiny{[GM$|$GM]}}
                \texttt{Pionós: níl (1,2) folamh, ach tá '─' ann.} \\
\\ 
\texttt{Tá 5 as 16 pionós carntha agaibh le chéile.} \\
\texttt{Déan iarracht arís, le do thoil!} \\
            }
        }
    }
    & & \\ \\

    \theutterance \stepcounter{utterance}  
    & & \multicolumn{4}{p{0.6\linewidth}}{
        \cellcolor[rgb]{0.9,0.9,0.9}{
            \makecell[{{p{\linewidth}}}]{
                \texttt{\tiny{[P1$\langle$GM]}}
                \texttt{Pionós: níl (1,2) folamh, ach tá '─' ann.} \\
\\ 
\texttt{Tá 5 as 16 pionós carntha agaibh le chéile.} \\
\texttt{Déan iarracht arís, le do thoil!} \\
            }
        }
    }
    & \\ \\

    \theutterance \stepcounter{utterance}  
    & \multicolumn{4}{p{0.6\linewidth}}{
        \cellcolor[rgb]{0.9,0.9,0.9}{
            \makecell[{{p{\linewidth}}}]{
                \texttt{\tiny{[P1$\rangle$GM]}}
                \texttt{BOG: M, (2,4)} \\
            }
        }
    }
    & & \\ \\

    \theutterance \stepcounter{utterance}  
    & & & \multicolumn{2}{p{0.3\linewidth}}{
        \cellcolor[rgb]{0.9,0.9,0.9}{
            \makecell[{{p{\linewidth}}}]{
                \texttt{\tiny{[GM$|$GM]}}
                \texttt{BOG: M, (2,4)} \\
            }
        }
    }
    & & \\ \\

    \theutterance \stepcounter{utterance}  
    & & & \multicolumn{2}{p{0.3\linewidth}}{
        \cellcolor[rgb]{0.9,0.9,0.9}{
            \makecell[{{p{\linewidth}}}]{
                \texttt{\tiny{[GM$|$GM]}}
                \texttt{Pionós: níl (2,4) folamh, ach tá '─' ann.} \\
\\ 
\texttt{Tá 6 as 16 pionós carntha agaibh le chéile.} \\
\texttt{Déan iarracht arís, le do thoil!} \\
            }
        }
    }
    & & \\ \\

    \theutterance \stepcounter{utterance}  
    & & \multicolumn{4}{p{0.6\linewidth}}{
        \cellcolor[rgb]{0.9,0.9,0.9}{
            \makecell[{{p{\linewidth}}}]{
                \texttt{\tiny{[P1$\langle$GM]}}
                \texttt{Pionós: níl (2,4) folamh, ach tá '─' ann.} \\
\\ 
\texttt{Tá 6 as 16 pionós carntha agaibh le chéile.} \\
\texttt{Déan iarracht arís, le do thoil!} \\
            }
        }
    }
    & \\ \\

    \theutterance \stepcounter{utterance}  
    & \multicolumn{4}{p{0.6\linewidth}}{
        \cellcolor[rgb]{0.9,0.9,0.9}{
            \makecell[{{p{\linewidth}}}]{
                \texttt{\tiny{[P1$\rangle$GM]}}
                \texttt{BOG: M, (3,2)} \\
            }
        }
    }
    & & \\ \\

    \theutterance \stepcounter{utterance}  
    & & & \multicolumn{2}{p{0.3\linewidth}}{
        \cellcolor[rgb]{0.9,0.9,0.9}{
            \makecell[{{p{\linewidth}}}]{
                \texttt{\tiny{[GM$|$GM]}}
                \texttt{BOG: M, (3,2)} \\
            }
        }
    }
    & & \\ \\

    \theutterance \stepcounter{utterance}  
    & & & \multicolumn{2}{p{0.3\linewidth}}{
        \cellcolor[rgb]{0.9,0.9,0.9}{
            \makecell[{{p{\linewidth}}}]{
                \texttt{\tiny{[GM$|$GM]}}
                \texttt{Pionós: níl (3,2) folamh, ach tá '┼' ann.} \\
\\ 
\texttt{Tá 7 as 16 pionós carntha agaibh le chéile.} \\
\texttt{Déan iarracht arís, le do thoil!} \\
            }
        }
    }
    & & \\ \\

    \theutterance \stepcounter{utterance}  
    & & \multicolumn{4}{p{0.6\linewidth}}{
        \cellcolor[rgb]{0.9,0.9,0.9}{
            \makecell[{{p{\linewidth}}}]{
                \texttt{\tiny{[P1$\langle$GM]}}
                \texttt{Pionós: níl (3,2) folamh, ach tá '┼' ann.} \\
\\ 
\texttt{Tá 7 as 16 pionós carntha agaibh le chéile.} \\
\texttt{Déan iarracht arís, le do thoil!} \\
            }
        }
    }
    & \\ \\

    \theutterance \stepcounter{utterance}  
    & \multicolumn{4}{p{0.6\linewidth}}{
        \cellcolor[rgb]{0.9,0.9,0.9}{
            \makecell[{{p{\linewidth}}}]{
                \texttt{\tiny{[P1$\rangle$GM]}}
                \texttt{BOG: M, (1,3)} \\
            }
        }
    }
    & & \\ \\

    \theutterance \stepcounter{utterance}  
    & & & \multicolumn{2}{p{0.3\linewidth}}{
        \cellcolor[rgb]{0.9,0.9,0.9}{
            \makecell[{{p{\linewidth}}}]{
                \texttt{\tiny{[GM$|$GM]}}
                \texttt{BOG: M, (1,3)} \\
            }
        }
    }
    & & \\ \\

    \theutterance \stepcounter{utterance}  
    & & & \multicolumn{2}{p{0.3\linewidth}}{
        \cellcolor[rgb]{0.9,0.9,0.9}{
            \makecell[{{p{\linewidth}}}]{
                \texttt{\tiny{[GM$|$GM]}}
                \texttt{M moved to (1, 3)} \\
            }
        }
    }
    & & \\ \\

    \theutterance \stepcounter{utterance}  
    & & \multicolumn{4}{p{0.6\linewidth}}{
        \cellcolor[rgb]{0.9,0.9,0.9}{
            \makecell[{{p{\linewidth}}}]{
                \texttt{\tiny{[GM$\rangle$P2]}}
                \texttt{Bogadh 'S' go (2,3) go rathúil. Seo mar atá do ghreille nuashonraithe:} \\
\texttt{```} \\
\texttt{ 1234567} \\
\texttt{╔══╤══╤═╗} \\
\texttt{║◌◌│◌◌│T║ 1} \\
\texttt{╟──┼──┴─╢ 2} \\
\texttt{║◌S│◌◌◌P║ 3} \\
\texttt{╟──┘◌◌MO║ 4} \\
\texttt{║D◌◌◌◌◌◌║ 5} \\
\texttt{║◌◌┌─┬──╢ 6} \\
\texttt{║◌◌│A│◌◌║ 7} \\
\texttt{╚══╧═╧══╝} \\
\texttt{```} \\
\texttt{Tá sibh ag imirt babhta 4 faoi láthair as uasmhéid 28.} \\
\texttt{Tá 7 as 16 pionós carntha agaibh le chéile.} \\
\texttt{Bhog an t{-}imreoir eile réad ar a ngreille.} \\
\\ 
\\ 
\texttt{Cad é do chéad ordú eile?} \\
            }
        }
    }
    & \\ \\

    \theutterance \stepcounter{utterance}  
    & & & \multicolumn{4}{p{0.6\linewidth}}{
        \cellcolor[rgb]{0.9,0.9,0.9}{
            \makecell[{{p{\linewidth}}}]{
                \texttt{\tiny{[GM$\langle$P2]}}
                \texttt{ABAIR: Tá 5 pionós carntha againn anois. Níor thuig mé an bhogadh a rinne an t{-}imreoir eile. An bhfuil S ag (2,3) ar do chuid freisin?} \\
            }
        }
    }
     \\ \\

    \theutterance \stepcounter{utterance}  
    & & & \multicolumn{2}{p{0.3\linewidth}}{
        \cellcolor[rgb]{0.9,0.9,0.9}{
            \makecell[{{p{\linewidth}}}]{
                \texttt{\tiny{[GM$|$GM]}}
                \texttt{ABAIR: Tá 5 pionós carntha againn anois. Níor thuig mé an bhogadh a rinne an t{-}imreoir eile. An bhfuil S ag (2,3) ar do chuid freisin?} \\
            }
        }
    }
    & & \\ \\

    \theutterance \stepcounter{utterance}  
    & & \multicolumn{4}{p{0.6\linewidth}}{
        \cellcolor[rgb]{0.9,0.9,0.9}{
            \makecell[{{p{\linewidth}}}]{
                \texttt{\tiny{[P1$\langle$GM]}}
                \texttt{Bogadh 'M' go (1,3) go rathúil. Seo mar atá do ghreille nuashonraithe:} \\
\texttt{```} \\
\texttt{ 1234567} \\
\texttt{╔══╤══╤═╗} \\
\texttt{║◌P│◌◌│◌║ 1} \\
\texttt{╟──┼──┴─╢ 2} \\
\texttt{║MS│◌◌◌O║ 3} \\
\texttt{╟──┘◌A◌◌║ 4} \\
\texttt{║D◌◌◌◌◌◌║ 5} \\
\texttt{║◌◌┌─┬──╢ 6} \\
\texttt{║T◌│◌│◌◌║ 7} \\
\texttt{╚══╧═╧══╝} \\
\texttt{```} \\
\texttt{Tá sibh ag imirt babhta 4 faoi láthair as uasmhéid 28.} \\
\texttt{Tá 7 as 16 pionós carntha agaibh le chéile.} \\
\texttt{Sheol an t{-}imreoir eile an teachtaireacht seo:} \\
\\ 
\\ 
\texttt{"Tá 5 pionós carntha againn anois. Níor thuig mé an bhogadh a rinne an t{-}imreoir eile. An bhfuil S ag (2,3) ar do chuid freisin?"} \\
\\ 
\\ 
\texttt{Cad é do chéad ordú eile?} \\
            }
        }
    }
    & \\ \\

    \theutterance \stepcounter{utterance}  
    & \multicolumn{4}{p{0.6\linewidth}}{
        \cellcolor[rgb]{0.9,0.9,0.9}{
            \makecell[{{p{\linewidth}}}]{
                \texttt{\tiny{[P1$\rangle$GM]}}
                \texttt{ABAIR: Tá S ag (2,3) agam anois. Táimid ag cur pionós ar fáil mar gheall ar bhogadh mícheart. Molaim go n{-}úsáidfimid an chéad bhabhta eile chun P a bhogadh go dtí (2,1) chun é a ailíniú leat, ach caithfidh M bogadh ar dtús. An bhfuil tú sásta go mbogaimid P go (2,1) anois? (Nóta: Tá M ag (1,3) agam anois, mar sin tá (2,1) folamh).} \\
            }
        }
    }
    & & \\ \\

    \theutterance \stepcounter{utterance}  
    & & & \multicolumn{2}{p{0.3\linewidth}}{
        \cellcolor[rgb]{0.9,0.9,0.9}{
            \makecell[{{p{\linewidth}}}]{
                \texttt{\tiny{[GM$|$GM]}}
                \texttt{ABAIR: Tá S ag (2,3) agam anois. Táimid ag cur pionós ar fáil mar gheall ar bhogadh mícheart. Molaim go n{-}úsáidfimid an chéad bhabhta eile chun P a bhogadh go dtí (2,1) chun é a ailíniú leat, ach caithfidh M bogadh ar dtús. An bhfuil tú sásta go mbogaimid P go (2,1) anois? (Nóta: Tá M ag (1,3) agam anois, mar sin tá (2,1) folamh).} \\
            }
        }
    }
    & & \\ \\

    \theutterance \stepcounter{utterance}  
    & & \multicolumn{4}{p{0.6\linewidth}}{
        \cellcolor[rgb]{0.9,0.9,0.9}{
            \makecell[{{p{\linewidth}}}]{
                \texttt{\tiny{[GM$\rangle$P2]}}
                \texttt{Cuireadh do theachtaireacht ar aghaidh chuig an imreoir eile.} \\
\\ 
\\ 
\texttt{Tá sibh ag imirt babhta 5 faoi láthair as uasmhéid 28.} \\
\texttt{Tá 7 as 16 pionós carntha agaibh le chéile.} \\
\texttt{Sheol an t{-}imreoir eile an teachtaireacht seo:} \\
\\ 
\\ 
\texttt{"Tá S ag (2,3) agam anois. Táimid ag cur pionós ar fáil mar gheall ar bhogadh mícheart. Molaim go n{-}úsáidfimid an chéad bhabhta eile chun P a bhogadh go dtí (2,1) chun é a ailíniú leat, ach caithfidh M bogadh ar dtús. An bhfuil tú sásta go mbogaimid P go (2,1) anois? (Nóta: Tá M ag (1,3) agam anois, mar sin tá (2,1) folamh)."} \\
\\ 
\\ 
\texttt{Cad é do chéad ordú eile?} \\
            }
        }
    }
    & \\ \\

    \theutterance \stepcounter{utterance}  
    & & & \multicolumn{4}{p{0.6\linewidth}}{
        \cellcolor[rgb]{0.9,0.9,0.9}{
            \makecell[{{p{\linewidth}}}]{
                \texttt{\tiny{[GM$\langle$P2]}}
                \texttt{ABAIR: Táim sásta. Bog P go (2,1).} \\
            }
        }
    }
     \\ \\

    \theutterance \stepcounter{utterance}  
    & & & \multicolumn{2}{p{0.3\linewidth}}{
        \cellcolor[rgb]{0.9,0.9,0.9}{
            \makecell[{{p{\linewidth}}}]{
                \texttt{\tiny{[GM$|$GM]}}
                \texttt{ABAIR: Táim sásta. Bog P go (2,1).} \\
            }
        }
    }
    & & \\ \\

    \theutterance \stepcounter{utterance}  
    & & \multicolumn{4}{p{0.6\linewidth}}{
        \cellcolor[rgb]{0.9,0.9,0.9}{
            \makecell[{{p{\linewidth}}}]{
                \texttt{\tiny{[P1$\langle$GM]}}
                \texttt{Cuireadh do theachtaireacht ar aghaidh chuig an imreoir eile.} \\
\\ 
\\ 
\texttt{Tá sibh ag imirt babhta 5 faoi láthair as uasmhéid 28.} \\
\texttt{Tá 7 as 16 pionós carntha agaibh le chéile.} \\
\texttt{Sheol an t{-}imreoir eile an teachtaireacht seo:} \\
\\ 
\\ 
\texttt{"Táim sásta. Bog P go (2,1)."} \\
\\ 
\\ 
\texttt{Cad é do chéad ordú eile?} \\
            }
        }
    }
    & \\ \\

    \theutterance \stepcounter{utterance}  
    & \multicolumn{4}{p{0.6\linewidth}}{
        \cellcolor[rgb]{0.9,0.9,0.9}{
            \makecell[{{p{\linewidth}}}]{
                \texttt{\tiny{[P1$\rangle$GM]}}
                \texttt{BOG: P, (2,1)} \\
            }
        }
    }
    & & \\ \\

    \theutterance \stepcounter{utterance}  
    & & & \multicolumn{2}{p{0.3\linewidth}}{
        \cellcolor[rgb]{0.9,0.9,0.9}{
            \makecell[{{p{\linewidth}}}]{
                \texttt{\tiny{[GM$|$GM]}}
                \texttt{BOG: P, (2,1)} \\
            }
        }
    }
    & & \\ \\

    \theutterance \stepcounter{utterance}  
    & & & \multicolumn{2}{p{0.3\linewidth}}{
        \cellcolor[rgb]{0.9,0.9,0.9}{
            \makecell[{{p{\linewidth}}}]{
                \texttt{\tiny{[GM$|$GM]}}
                \texttt{Pionós: níl (2,1) folamh, ach tá 'P' ann.} \\
\\ 
\texttt{Tá 8 as 16 pionós carntha agaibh le chéile.} \\
\texttt{Déan iarracht arís, le do thoil!} \\
            }
        }
    }
    & & \\ \\

    \theutterance \stepcounter{utterance}  
    & & \multicolumn{4}{p{0.6\linewidth}}{
        \cellcolor[rgb]{0.9,0.9,0.9}{
            \makecell[{{p{\linewidth}}}]{
                \texttt{\tiny{[P1$\langle$GM]}}
                \texttt{Pionós: níl (2,1) folamh, ach tá 'P' ann.} \\
\\ 
\texttt{Tá 8 as 16 pionós carntha agaibh le chéile.} \\
\texttt{Déan iarracht arís, le do thoil!} \\
            }
        }
    }
    & \\ \\

    \theutterance \stepcounter{utterance}  
    & \multicolumn{4}{p{0.6\linewidth}}{
        \cellcolor[rgb]{0.9,0.9,0.9}{
            \makecell[{{p{\linewidth}}}]{
                \texttt{\tiny{[P1$\rangle$GM]}}
                \texttt{ABAIR: Tá botún déanta agam. Tá P ag (1,1) fós. Tá M ag (1,3). Tá S ag (2,3). Tá O ag (4,3). Tá A ag (3,4). Tá D ag (1,5). Tá T ag (1,7).} \\
\texttt{Táimid ag iarraidh P a bhogadh go (2,1). Ach tá P ag (1,1).} \\
\texttt{Molaim go mbogaimid P ó (1,1) go (2,1) anois. Ach tá P ag (1,1) agus tá (2,1) folamh?} \\
\texttt{Seiceáil:} \\
\texttt{Ró 1: P ag (1,1).} \\
\texttt{Ró 2: M ag (1,3), S ag (2,3).} \\
\texttt{Ró 3: O ag (4,3).} \\
\texttt{Ró 4: A ag (3,4).} \\
\texttt{Ró 5: D ag (1,5).} \\
\texttt{Ró 7: T ag (1,7).} \\
\\ 
\texttt{Tá (2,1) folamh?} \\
\texttt{Ón gcéad ghreille:} \\
\texttt{(1,1) P} \\
\texttt{(2,1) M} \\
\texttt{(3,1) ◌} \\
\texttt{(4,1) ◌} \\
\texttt{(5,1) ◌} \\
\texttt{(6,1) ◌} \\
\texttt{(7,1) ◌} \\
\\ 
\texttt{Tar éis M a bhogadh go (1,3):} \\
\texttt{(1,1) P} \\
\texttt{(1,3) M} \\
\texttt{(2,3) S} \\
\texttt{(4,3) O} \\
\texttt{(3,4) A} \\
\texttt{(1,5) D} \\
\texttt{(1,7) T} \\
\\ 
\texttt{Tá (2,1) folamh.} \\
\texttt{Mar sin, is féidir liom P a bhogadh ó (1,1) go (2,1).} \\
\\ 
\texttt{BOG: P, (2,1)} \\
            }
        }
    }
    & & \\ \\

    \theutterance \stepcounter{utterance}  
    & & & \multicolumn{2}{p{0.3\linewidth}}{
        \cellcolor[rgb]{0.9,0.9,0.9}{
            \makecell[{{p{\linewidth}}}]{
                \texttt{\tiny{[GM$|$GM]}}
                \texttt{ABAIR: Tá botún déanta agam. Tá P ag (1,1) fós. Tá M ag (1,3). Tá S ag (2,3). Tá O ag (4,3). Tá A ag (3,4). Tá D ag (1,5). Tá T ag (1,7).} \\
\texttt{Táimid ag iarraidh P a bhogadh go (2,1). Ach tá P ag (1,1).} \\
\texttt{Molaim go mbogaimid P ó (1,1) go (2,1) anois. Ach tá P ag (1,1) agus tá (2,1) folamh?} \\
\texttt{Seiceáil:} \\
\texttt{Ró 1: P ag (1,1).} \\
\texttt{Ró 2: M ag (1,3), S ag (2,3).} \\
\texttt{Ró 3: O ag (4,3).} \\
\texttt{Ró 4: A ag (3,4).} \\
\texttt{Ró 5: D ag (1,5).} \\
\texttt{Ró 7: T ag (1,7).} \\
\\ 
\texttt{Tá (2,1) folamh?} \\
\texttt{Ón gcéad ghreille:} \\
\texttt{(1,1) P} \\
\texttt{(2,1) M} \\
\texttt{(3,1) ◌} \\
\texttt{(4,1) ◌} \\
\texttt{(5,1) ◌} \\
\texttt{(6,1) ◌} \\
\texttt{(7,1) ◌} \\
\\ 
\texttt{Tar éis M a bhogadh go (1,3):} \\
\texttt{(1,1) P} \\
\texttt{(1,3) M} \\
\texttt{(2,3) S} \\
\texttt{(4,3) O} \\
\texttt{(3,4) A} \\
\texttt{(1,5) D} \\
\texttt{(1,7) T} \\
\\ 
\texttt{Tá (2,1) folamh.} \\
\texttt{Mar sin, is féidir liom P a bhogadh ó (1,1) go (2,1).} \\
\\ 
\texttt{BOG: P, (2,1)} \\
            }
        }
    }
    & & \\ \\

    \theutterance \stepcounter{utterance}  
    & & & \multicolumn{2}{p{0.3\linewidth}}{
        \cellcolor[rgb]{0.9,0.9,0.9}{
            \makecell[{{p{\linewidth}}}]{
                \texttt{\tiny{[GM$|$GM]}}
                \texttt{Invalid response format: ABAIR: Tá botún déanta agam. Tá P ag (1,1) fós. Tá M ag (1,3). Tá S ag (2,3). Tá O ag (4,3). Tá A ag (3,4). Tá D ag (1,5). Tá T ag (1,7).} \\
\texttt{Táimid ag iarraidh P a bhogadh go (2,1). Ach tá P ag (1,1).} \\
\texttt{Molaim go mbogaimid P ó (1,1) go (2,1) anois. Ach tá P ag (1,1) agus tá (2,1) folamh?} \\
\texttt{Seiceáil:} \\
\texttt{Ró 1: P ag (1,1).} \\
\texttt{Ró 2: M ag (1,3), S ag (2,3).} \\
\texttt{Ró 3: O ag (4,3).} \\
\texttt{Ró 4: A ag (3,4).} \\
\texttt{Ró 5: D ag (1,5).} \\
\texttt{Ró 7: T ag (1,7).} \\
\\ 
\texttt{Tá (2,1) folamh?} \\
\texttt{Ón gcéad ghreille:} \\
\texttt{(1,1) P} \\
\texttt{(2,1) M} \\
\texttt{(3,1) ◌} \\
\texttt{(4,1) ◌} \\
\texttt{(5,1) ◌} \\
\texttt{(6,1) ◌} \\
\texttt{(7,1) ◌} \\
\\ 
\texttt{Tar éis M a bhogadh go (1,3):} \\
\texttt{(1,1) P} \\
\texttt{(1,3) M} \\
\texttt{(2,3) S} \\
\texttt{(4,3) O} \\
\texttt{(3,4) A} \\
\texttt{(1,5) D} \\
\texttt{(1,7) T} \\
\\ 
\texttt{Tá (2,1) folamh.} \\
\texttt{Mar sin, is féidir liom P a bhogadh ó (1,1) go (2,1).} \\
\\ 
\texttt{BOG: P, (2,1)} \\
            }
        }
    }
    & & \\ \\

    \theutterance \stepcounter{utterance}  
    & & \multicolumn{4}{p{0.6\linewidth}}{
        \cellcolor[rgb]{0.9,0.9,0.9}{
            \makecell[{{p{\linewidth}}}]{
                \texttt{\tiny{[P1$\langle$GM]}}
                \texttt{Pionós: Tá níos mó ná ordú amháin i do theachtaireacht!} \\
\\ 
\texttt{Cinntigh nach bhfuil i do fhreagra ach `ABAIR: <TEACHTAIREACHT>` nó `BOG: <RUD>, (<X>, <Y>)`, agus tada eile!} \\
\\ 
\texttt{Tá 9 as 16 pionós carntha agaibh le chéile.} \\
\texttt{Déan iarracht arís, le do thoil!} \\
            }
        }
    }
    & \\ \\

    \theutterance \stepcounter{utterance}  
    & \multicolumn{4}{p{0.6\linewidth}}{
        \cellcolor[rgb]{0.9,0.9,0.9}{
            \makecell[{{p{\linewidth}}}]{
                \texttt{\tiny{[P1$\rangle$GM]}}
                \texttt{BOG: P, (2,1)} \\
            }
        }
    }
    & & \\ \\

    \theutterance \stepcounter{utterance}  
    & & & \multicolumn{2}{p{0.3\linewidth}}{
        \cellcolor[rgb]{0.9,0.9,0.9}{
            \makecell[{{p{\linewidth}}}]{
                \texttt{\tiny{[GM$|$GM]}}
                \texttt{BOG: P, (2,1)} \\
            }
        }
    }
    & & \\ \\

    \theutterance \stepcounter{utterance}  
    & & & \multicolumn{2}{p{0.3\linewidth}}{
        \cellcolor[rgb]{0.9,0.9,0.9}{
            \makecell[{{p{\linewidth}}}]{
                \texttt{\tiny{[GM$|$GM]}}
                \texttt{Pionós: níl (2,1) folamh, ach tá 'P' ann.} \\
\\ 
\texttt{Tá 10 as 16 pionós carntha agaibh le chéile.} \\
\texttt{Déan iarracht arís, le do thoil!} \\
            }
        }
    }
    & & \\ \\

    \theutterance \stepcounter{utterance}  
    & & \multicolumn{4}{p{0.6\linewidth}}{
        \cellcolor[rgb]{0.9,0.9,0.9}{
            \makecell[{{p{\linewidth}}}]{
                \texttt{\tiny{[P1$\langle$GM]}}
                \texttt{Pionós: níl (2,1) folamh, ach tá 'P' ann.} \\
\\ 
\texttt{Tá 10 as 16 pionós carntha agaibh le chéile.} \\
\texttt{Déan iarracht arís, le do thoil!} \\
            }
        }
    }
    & \\ \\

    \theutterance \stepcounter{utterance}  
    & \multicolumn{4}{p{0.6\linewidth}}{
        \cellcolor[rgb]{0.9,0.9,0.9}{
            \makecell[{{p{\linewidth}}}]{
                \texttt{\tiny{[P1$\rangle$GM]}}
                \texttt{ABAIR: Tá botún agam. Ní féidir liom P a bhogadh go (2,1) toisc go bhfuil P ann cheana féin ag (1,1). Táimid ag iarraidh P a ailíniú leat. Tá P agat ag (5,3). Tá P agam ag (1,1).} \\
\texttt{Molaim go n{-}úsáidfimid suíomh eile. An bhfuil tú sásta go mbogaimid P go (1,2) ar dtús? (Seiceáil: Tá (1,2) folamh? Ní hea, tá '{-}' ann. Tá (1,2) folamh? Ní hea, tá '{-}' ann. Tá (1,2) folamh? Ní hea, tá '{-}' ann.} \\
\texttt{Táimid ag déanamh botúin arís.} \\
\texttt{Táimid ag iarraidh P a bhogadh go dtí suíomh atá folamh.} \\
\texttt{Tá P ag (1,1).} \\
\texttt{Tá (1,2) folamh? Ní hea, tá '{-}' ann.} \\
\texttt{Tá (2,1) folamh? Ní hea, tá P ann (tar éis mo iarracht theipthe).} \\
\texttt{Tá (1,0) lasmuigh.} \\
\texttt{Tá (0,1) lasmuigh.} \\
\\ 
\texttt{Ní mór dúinn P a bhogadh go dtí suíomh atá folamh.} \\
\texttt{Tá P ag (1,1).} \\
\texttt{Tá (2,1) folamh? Ní hea, tá P ann.} \\
\texttt{Tá (1,2) folamh? Ní hea, tá '{-}' ann.} \\
\texttt{Tá (0,1) lasmuigh.} \\
\texttt{Tá (1,0) lasmuigh.} \\
\\ 
\texttt{Níl aon bhealach eile agam P a bhogadh?} \\
\texttt{Tá P ag (1,1).} \\
\texttt{Tá (2,1) folamh? Ní hea, tá P ann.} \\
\texttt{Tá (1,2) folamh? Ní hea, tá '{-}' ann.} \\
\\ 
\texttt{Ní mór dúinn P a bhogadh go dtí suíomh atá folamh.} \\
\texttt{Tá P ag (1,1).} \\
\texttt{Tá (2,1) folamh? Ní hea, tá P ann.} \\
\texttt{Tá (1,2) folamh? Ní hea, tá '{-}' ann.} \\
\\ 
\texttt{Níl aon bhealach eile agam P a bhogadh?} \\
\texttt{Tá P ag (1,1).} \\
\texttt{Tá (2,1) folamh? Ní hea, tá P ann.} \\
\texttt{Tá (1,2) folamh? Ní hea, tá '{-}' ann.} \\
\\ 
\texttt{Ní mór dúinn P a bhogadh go dtí suíomh atá folamh.} \\
\texttt{Tá P ag (1,1).} \\
\texttt{Tá (2,1) folamh? Ní hea, tá P ann.} \\
\texttt{Tá (1,2) folamh? Ní hea, tá '{-}' ann.} \\
\\ 
\texttt{Níl aon bhealach eile agam P a bhogadh?} \\
\texttt{Tá P ag (1,1).} \\
\texttt{Tá (2,1) folamh? Ní hea, tá P ann.} \\
\texttt{Tá (1,2) folamh? Ní hea, tá '{-}' ann.} \\
\\ 
\texttt{Ní mór dúinn P a bhogadh go dtí suíomh atá folamh.} \\
\texttt{Tá P ag (1,1).} \\
\texttt{Tá (2,1) folamh? Ní hea, tá P ann.} \\
\texttt{Tá (1,2) folamh? Ní hea, tá '{-}' ann.} \\
\\ 
\texttt{Níl aon bhealach eile agam P a bhogadh?} \\
\texttt{Tá P ag (1,1).} \\
\texttt{Tá (2,1) folamh? Ní hea, tá P ann.} \\
\texttt{Tá (1,2) folamh? Ní hea, tá '{-}' ann.} \\
\\ 
\texttt{Ní mór dúinn P a bhogadh go dtí suíomh atá folamh.} \\
\texttt{Tá P ag (1,1).} \\
\texttt{Tá (2,1) folamh? Ní hea, tá P ann.} \\
\texttt{Tá (1,2) folamh? Ní hea, tá '{-}' ann.} \\
\\ 
\texttt{Níl aon bhealach eile agam P a bhogadh?} \\
\texttt{Tá P ag (1,1).} \\
\texttt{Tá (2,1) folamh? Ní hea, tá P ann.} \\
\texttt{Tá (1,2) folamh? Ní hea, tá '{-}' ann.} \\
\\ 
\texttt{Ní mór dúinn P a bhogadh go dtí suíomh atá folamh.} \\
\texttt{Tá P ag (1,1).} \\
\texttt{Tá (2,1) folamh? Ní hea, tá P ann.} \\
\texttt{Tá (1,2) folamh? Ní hea, tá '{-}' ann.} \\
\\ 
\texttt{Níl aon bhealach eile agam P a bhogadh?} \\
\texttt{Tá P ag (1,1).} \\
\texttt{Tá (2,1) folamh? Ní hea, tá P ann.} \\
\texttt{Tá (1,2) folamh? Ní hea, tá '{-}' ann.} \\
\\ 
\texttt{Ní mór dúinn P a bhogadh go dtí suíomh atá folamh.} \\
\texttt{Tá P ag (1,1).} \\
\texttt{Tá (2,1) folamh? Ní hea, tá P ann.} \\
\texttt{Tá (1,2) folamh? Ní hea, tá '{-}' ann.} \\
\\ 
\texttt{Níl aon bhealach eile agam P a bhogadh?} \\
\texttt{Tá P ag (1,1).} \\
\texttt{Tá (2,1) folamh? Ní hea, tá P ann.} \\
\texttt{Tá (1,2) folamh? Ní hea, tá '{-}' ann.} \\
\\ 
\texttt{Ní mór dúinn P a bhogadh go dtí suíomh atá folamh.} \\
\texttt{Tá P ag (1,1).} \\
\texttt{Tá (2,1) folamh? Ní hea, tá P ann.} \\
\texttt{Tá (1,2) folamh? Ní hea, tá '{-}' ann.} \\
\\ 
\texttt{Níl aon bhealach eile agam P a bhogadh?} \\
\texttt{Tá P ag (1,1).} \\
\texttt{Tá (2,1) folamh? Ní hea, tá P ann.} \\
\texttt{Tá (1,2) folamh? Ní hea, tá '{-}' ann.} \\
\\ 
\texttt{Ní mór dúinn P a bhogadh go dtí suíomh atá folamh.} \\
\texttt{Tá P ag (1,1).} \\
\texttt{Tá (2,1) folamh? Ní hea, tá P ann.} \\
\texttt{Tá (1,2) folamh? Ní hea, tá '{-}' ann.} \\
\\ 
\texttt{Níl aon bhealach eile agam P a bhogadh?} \\
\texttt{Tá P ag (1,1).} \\
\texttt{Tá (2,1) folamh? Ní hea, tá P ann.} \\
\texttt{Tá (1,2) folamh? Ní hea, tá '{-}' ann.} \\
\\ 
\texttt{Ní mór dúinn P a bhogadh go dtí suíomh atá folamh.} \\
\texttt{Tá P ag (1,1).} \\
\texttt{Tá (2,1) folamh? Ní hea, tá P ann.} \\
\texttt{Tá (1,2) folamh? Ní hea, tá '{-}' ann.} \\
\\ 
\texttt{Níl aon bhealach eile agam P a bhogadh?} \\
\texttt{Tá P ag (1,1).} \\
\texttt{Tá (2,1) folamh? Ní hea, tá P ann.} \\
\texttt{Tá (1,2) folamh? Ní hea, tá '{-}' ann.} \\
\\ 
\texttt{Ní mór dúinn P a bhogadh go dtí suíomh atá folamh.} \\
\texttt{Tá P ag (1,1).} \\
\texttt{Tá (2,1) folamh? Ní hea, tá P ann.} \\
\texttt{Tá (1,2) folamh? Ní hea, tá '{-}' ann.} \\
\\ 
\texttt{Níl aon bhealach eile agam P a bhogadh?} \\
\texttt{Tá P ag (1,1).} \\
\texttt{Tá (2,1) folamh? Ní hea, tá P ann.} \\
\texttt{Tá (1,2) folamh? Ní hea, tá '{-}' ann.} \\
\\ 
\texttt{Ní mór dúinn P a bhogadh go dtí suíomh atá folamh.} \\
\texttt{Tá P ag (1,1).} \\
\texttt{Tá (2,1) folamh? N} \\
            }
        }
    }
    & & \\ \\

    \theutterance \stepcounter{utterance}  
    & & & \multicolumn{2}{p{0.3\linewidth}}{
        \cellcolor[rgb]{0.9,0.9,0.9}{
            \makecell[{{p{\linewidth}}}]{
                \texttt{\tiny{[GM$|$GM]}}
                \texttt{ABAIR: Tá botún agam. Ní féidir liom P a bhogadh go (2,1) toisc go bhfuil P ann cheana féin ag (1,1). Táimid ag iarraidh P a ailíniú leat. Tá P agat ag (5,3). Tá P agam ag (1,1).} \\
\texttt{Molaim go n{-}úsáidfimid suíomh eile. An bhfuil tú sásta go mbogaimid P go (1,2) ar dtús? (Seiceáil: Tá (1,2) folamh? Ní hea, tá '{-}' ann. Tá (1,2) folamh? Ní hea, tá '{-}' ann. Tá (1,2) folamh? Ní hea, tá '{-}' ann.} \\
\texttt{Táimid ag déanamh botúin arís.} \\
\texttt{Táimid ag iarraidh P a bhogadh go dtí suíomh atá folamh.} \\
\texttt{Tá P ag (1,1).} \\
\texttt{Tá (1,2) folamh? Ní hea, tá '{-}' ann.} \\
\texttt{Tá (2,1) folamh? Ní hea, tá P ann (tar éis mo iarracht theipthe).} \\
\texttt{Tá (1,0) lasmuigh.} \\
\texttt{Tá (0,1) lasmuigh.} \\
\\ 
\texttt{Ní mór dúinn P a bhogadh go dtí suíomh atá folamh.} \\
\texttt{Tá P ag (1,1).} \\
\texttt{Tá (2,1) folamh? Ní hea, tá P ann.} \\
\texttt{Tá (1,2) folamh? Ní hea, tá '{-}' ann.} \\
\texttt{Tá (0,1) lasmuigh.} \\
\texttt{Tá (1,0) lasmuigh.} \\
\\ 
\texttt{Níl aon bhealach eile agam P a bhogadh?} \\
\texttt{Tá P ag (1,1).} \\
\texttt{Tá (2,1) folamh? Ní hea, tá P ann.} \\
\texttt{Tá (1,2) folamh? Ní hea, tá '{-}' ann.} \\
\\ 
\texttt{Ní mór dúinn P a bhogadh go dtí suíomh atá folamh.} \\
\texttt{Tá P ag (1,1).} \\
\texttt{Tá (2,1) folamh? Ní hea, tá P ann.} \\
\texttt{Tá (1,2) folamh? Ní hea, tá '{-}' ann.} \\
\\ 
\texttt{Níl aon bhealach eile agam P a bhogadh?} \\
\texttt{Tá P ag (1,1).} \\
\texttt{Tá (2,1) folamh? Ní hea, tá P ann.} \\
\texttt{Tá (1,2) folamh? Ní hea, tá '{-}' ann.} \\
\\ 
\texttt{Ní mór dúinn P a bhogadh go dtí suíomh atá folamh.} \\
\texttt{Tá P ag (1,1).} \\
\texttt{Tá (2,1) folamh? Ní hea, tá P ann.} \\
\texttt{Tá (1,2) folamh? Ní hea, tá '{-}' ann.} \\
\\ 
\texttt{Níl aon bhealach eile agam P a bhogadh?} \\
\texttt{Tá P ag (1,1).} \\
\texttt{Tá (2,1) folamh? Ní hea, tá P ann.} \\
\texttt{Tá (1,2) folamh? Ní hea, tá '{-}' ann.} \\
\\ 
\texttt{Ní mór dúinn P a bhogadh go dtí suíomh atá folamh.} \\
\texttt{Tá P ag (1,1).} \\
\texttt{Tá (2,1) folamh? Ní hea, tá P ann.} \\
\texttt{Tá (1,2) folamh? Ní hea, tá '{-}' ann.} \\
\\ 
\texttt{Níl aon bhealach eile agam P a bhogadh?} \\
\texttt{Tá P ag (1,1).} \\
\texttt{Tá (2,1) folamh? Ní hea, tá P ann.} \\
\texttt{Tá (1,2) folamh? Ní hea, tá '{-}' ann.} \\
\\ 
\texttt{Ní mór dúinn P a bhogadh go dtí suíomh atá folamh.} \\
\texttt{Tá P ag (1,1).} \\
\texttt{Tá (2,1) folamh? Ní hea, tá P ann.} \\
\texttt{Tá (1,2) folamh? Ní hea, tá '{-}' ann.} \\
\\ 
\texttt{Níl aon bhealach eile agam P a bhogadh?} \\
\texttt{Tá P ag (1,1).} \\
\texttt{Tá (2,1) folamh? Ní hea, tá P ann.} \\
\texttt{Tá (1,2) folamh? Ní hea, tá '{-}' ann.} \\
\\ 
\texttt{Ní mór dúinn P a bhogadh go dtí suíomh atá folamh.} \\
\texttt{Tá P ag (1,1).} \\
\texttt{Tá (2,1) folamh? Ní hea, tá P ann.} \\
\texttt{Tá (1,2) folamh? Ní hea, tá '{-}' ann.} \\
\\ 
\texttt{Níl aon bhealach eile agam P a bhogadh?} \\
\texttt{Tá P ag (1,1).} \\
\texttt{Tá (2,1) folamh? Ní hea, tá P ann.} \\
\texttt{Tá (1,2) folamh? Ní hea, tá '{-}' ann.} \\
\\ 
\texttt{Ní mór dúinn P a bhogadh go dtí suíomh atá folamh.} \\
\texttt{Tá P ag (1,1).} \\
\texttt{Tá (2,1) folamh? Ní hea, tá P ann.} \\
\texttt{Tá (1,2) folamh? Ní hea, tá '{-}' ann.} \\
\\ 
\texttt{Níl aon bhealach eile agam P a bhogadh?} \\
\texttt{Tá P ag (1,1).} \\
\texttt{Tá (2,1) folamh? Ní hea, tá P ann.} \\
\texttt{Tá (1,2) folamh? Ní hea, tá '{-}' ann.} \\
\\ 
\texttt{Ní mór dúinn P a bhogadh go dtí suíomh atá folamh.} \\
\texttt{Tá P ag (1,1).} \\
\texttt{Tá (2,1) folamh? Ní hea, tá P ann.} \\
\texttt{Tá (1,2) folamh? Ní hea, tá '{-}' ann.} \\
\\ 
\texttt{Níl aon bhealach eile agam P a bhogadh?} \\
\texttt{Tá P ag (1,1).} \\
\texttt{Tá (2,1) folamh? Ní hea, tá P ann.} \\
\texttt{Tá (1,2) folamh? Ní hea, tá '{-}' ann.} \\
\\ 
\texttt{Ní mór dúinn P a bhogadh go dtí suíomh atá folamh.} \\
\texttt{Tá P ag (1,1).} \\
\texttt{Tá (2,1) folamh? Ní hea, tá P ann.} \\
\texttt{Tá (1,2) folamh? Ní hea, tá '{-}' ann.} \\
\\ 
\texttt{Níl aon bhealach eile agam P a bhogadh?} \\
\texttt{Tá P ag (1,1).} \\
\texttt{Tá (2,1) folamh? Ní hea, tá P ann.} \\
\texttt{Tá (1,2) folamh? Ní hea, tá '{-}' ann.} \\
\\ 
\texttt{Ní mór dúinn P a bhogadh go dtí suíomh atá folamh.} \\
\texttt{Tá P ag (1,1).} \\
\texttt{Tá (2,1) folamh? Ní hea, tá P ann.} \\
\texttt{Tá (1,2) folamh? Ní hea, tá '{-}' ann.} \\
\\ 
\texttt{Níl aon bhealach eile agam P a bhogadh?} \\
\texttt{Tá P ag (1,1).} \\
\texttt{Tá (2,1) folamh? Ní hea, tá P ann.} \\
\texttt{Tá (1,2) folamh? Ní hea, tá '{-}' ann.} \\
\\ 
\texttt{Ní mór dúinn P a bhogadh go dtí suíomh atá folamh.} \\
\texttt{Tá P ag (1,1).} \\
\texttt{Tá (2,1) folamh? Ní hea, tá P ann.} \\
\texttt{Tá (1,2) folamh? Ní hea, tá '{-}' ann.} \\
\\ 
\texttt{Níl aon bhealach eile agam P a bhogadh?} \\
\texttt{Tá P ag (1,1).} \\
\texttt{Tá (2,1) folamh? Ní hea, tá P ann.} \\
\texttt{Tá (1,2) folamh? Ní hea, tá '{-}' ann.} \\
\\ 
\texttt{Ní mór dúinn P a bhogadh go dtí suíomh atá folamh.} \\
\texttt{Tá P ag (1,1).} \\
\texttt{Tá (2,1) folamh? N} \\
            }
        }
    }
    & & \\ \\

    \theutterance \stepcounter{utterance}  
    & & & \multicolumn{2}{p{0.3\linewidth}}{
        \cellcolor[rgb]{0.9,0.9,0.9}{
            \makecell[{{p{\linewidth}}}]{
                \texttt{\tiny{[GM$|$GM]}}
                \texttt{Invalid format: tail is not empty: '} \\
\texttt{Molaim go n{-}úsáidfimid suíomh eile. An bhfuil tú sásta go mbogaimid P go (1,2) ar dtús? (Seiceáil: Tá (1,2) folamh? Ní hea, tá '{-}' ann. Tá (1,2) folamh? Ní hea, tá '{-}' ann. Tá (1,2) folamh? Ní hea, tá '{-}' ann.} \\
\texttt{Táimid ag déanamh botúin arís.} \\
\texttt{Táimid ag iarraidh P a bhogadh go dtí suíomh atá folamh.} \\
\texttt{Tá P ag (1,1).} \\
\texttt{Tá (1,2) folamh? Ní hea, tá '{-}' ann.} \\
\texttt{Tá (2,1) folamh? Ní hea, tá P ann (tar éis mo iarracht theipthe).} \\
\texttt{Tá (1,0) lasmuigh.} \\
\texttt{Tá (0,1) lasmuigh.} \\
\\ 
\texttt{Ní mór dúinn P a bhogadh go dtí suíomh atá folamh.} \\
\texttt{Tá P ag (1,1).} \\
\texttt{Tá (2,1) folamh? Ní hea, tá P ann.} \\
\texttt{Tá (1,2) folamh? Ní hea, tá '{-}' ann.} \\
\texttt{Tá (0,1) lasmuigh.} \\
\texttt{Tá (1,0) lasmuigh.} \\
\\ 
\texttt{Níl aon bhealach eile agam P a bhogadh?} \\
\texttt{Tá P ag (1,1).} \\
\texttt{Tá (2,1) folamh? Ní hea, tá P ann.} \\
\texttt{Tá (1,2) folamh? Ní hea, tá '{-}' ann.} \\
\\ 
\texttt{Ní mór dúinn P a bhogadh go dtí suíomh atá folamh.} \\
\texttt{Tá P ag (1,1).} \\
\texttt{Tá (2,1) folamh? Ní hea, tá P ann.} \\
\texttt{Tá (1,2) folamh? Ní hea, tá '{-}' ann.} \\
\\ 
\texttt{Níl aon bhealach eile agam P a bhogadh?} \\
\texttt{Tá P ag (1,1).} \\
\texttt{Tá (2,1) folamh? Ní hea, tá P ann.} \\
\texttt{Tá (1,2) folamh? Ní hea, tá '{-}' ann.} \\
\\ 
\texttt{Ní mór dúinn P a bhogadh go dtí suíomh atá folamh.} \\
\texttt{Tá P ag (1,1).} \\
\texttt{Tá (2,1) folamh? Ní hea, tá P ann.} \\
\texttt{Tá (1,2) folamh? Ní hea, tá '{-}' ann.} \\
\\ 
\texttt{Níl aon bhealach eile agam P a bhogadh?} \\
\texttt{Tá P ag (1,1).} \\
\texttt{Tá (2,1) folamh? Ní hea, tá P ann.} \\
\texttt{Tá (1,2) folamh? Ní hea, tá '{-}' ann.} \\
\\ 
\texttt{Ní mór dúinn P a bhogadh go dtí suíomh atá folamh.} \\
\texttt{Tá P ag (1,1).} \\
\texttt{Tá (2,1) folamh? Ní hea, tá P ann.} \\
\texttt{Tá (1,2) folamh? Ní hea, tá '{-}' ann.} \\
\\ 
\texttt{Níl aon bhealach eile agam P a bhogadh?} \\
\texttt{Tá P ag (1,1).} \\
\texttt{Tá (2,1) folamh? Ní hea, tá P ann.} \\
\texttt{Tá (1,2) folamh? Ní hea, tá '{-}' ann.} \\
\\ 
\texttt{Ní mór dúinn P a bhogadh go dtí suíomh atá folamh.} \\
\texttt{Tá P ag (1,1).} \\
\texttt{Tá (2,1) folamh? Ní hea, tá P ann.} \\
\texttt{Tá (1,2) folamh? Ní hea, tá '{-}' ann.} \\
\\ 
\texttt{Níl aon bhealach eile agam P a bhogadh?} \\
\texttt{Tá P ag (1,1).} \\
\texttt{Tá (2,1) folamh? Ní hea, tá P ann.} \\
\texttt{Tá (1,2) folamh? Ní hea, tá '{-}' ann.} \\
\\ 
\texttt{Ní mór dúinn P a bhogadh go dtí suíomh atá folamh.} \\
\texttt{Tá P ag (1,1).} \\
\texttt{Tá (2,1) folamh? Ní hea, tá P ann.} \\
\texttt{Tá (1,2) folamh? Ní hea, tá '{-}' ann.} \\
\\ 
\texttt{Níl aon bhealach eile agam P a bhogadh?} \\
\texttt{Tá P ag (1,1).} \\
\texttt{Tá (2,1) folamh? Ní hea, tá P ann.} \\
\texttt{Tá (1,2) folamh? Ní hea, tá '{-}' ann.} \\
\\ 
\texttt{Ní mór dúinn P a bhogadh go dtí suíomh atá folamh.} \\
\texttt{Tá P ag (1,1).} \\
\texttt{Tá (2,1) folamh? Ní hea, tá P ann.} \\
\texttt{Tá (1,2) folamh? Ní hea, tá '{-}' ann.} \\
\\ 
\texttt{Níl aon bhealach eile agam P a bhogadh?} \\
\texttt{Tá P ag (1,1).} \\
\texttt{Tá (2,1) folamh? Ní hea, tá P ann.} \\
\texttt{Tá (1,2) folamh? Ní hea, tá '{-}' ann.} \\
\\ 
\texttt{Ní mór dúinn P a bhogadh go dtí suíomh atá folamh.} \\
\texttt{Tá P ag (1,1).} \\
\texttt{Tá (2,1) folamh? Ní hea, tá P ann.} \\
\texttt{Tá (1,2) folamh? Ní hea, tá '{-}' ann.} \\
\\ 
\texttt{Níl aon bhealach eile agam P a bhogadh?} \\
\texttt{Tá P ag (1,1).} \\
\texttt{Tá (2,1) folamh? Ní hea, tá P ann.} \\
\texttt{Tá (1,2) folamh? Ní hea, tá '{-}' ann.} \\
\\ 
\texttt{Ní mór dúinn P a bhogadh go dtí suíomh atá folamh.} \\
\texttt{Tá P ag (1,1).} \\
\texttt{Tá (2,1) folamh? Ní hea, tá P ann.} \\
\texttt{Tá (1,2) folamh? Ní hea, tá '{-}' ann.} \\
\\ 
\texttt{Níl aon bhealach eile agam P a bhogadh?} \\
\texttt{Tá P ag (1,1).} \\
\texttt{Tá (2,1) folamh? Ní hea, tá P ann.} \\
\texttt{Tá (1,2) folamh? Ní hea, tá '{-}' ann.} \\
\\ 
\texttt{Ní mór dúinn P a bhogadh go dtí suíomh atá folamh.} \\
\texttt{Tá P ag (1,1).} \\
\texttt{Tá (2,1) folamh? Ní hea, tá P ann.} \\
\texttt{Tá (1,2) folamh? Ní hea, tá '{-}' ann.} \\
\\ 
\texttt{Níl aon bhealach eile agam P a bhogadh?} \\
\texttt{Tá P ag (1,1).} \\
\texttt{Tá (2,1) folamh? Ní hea, tá P ann.} \\
\texttt{Tá (1,2) folamh? Ní hea, tá '{-}' ann.} \\
\\ 
\texttt{Ní mór dúinn P a bhogadh go dtí suíomh atá folamh.} \\
\texttt{Tá P ag (1,1).} \\
\texttt{Tá (2,1) folamh? Ní hea, tá P ann.} \\
\texttt{Tá (1,2) folamh? Ní hea, tá '{-}' ann.} \\
\\ 
\texttt{Níl aon bhealach eile agam P a bhogadh?} \\
\texttt{Tá P ag (1,1).} \\
\texttt{Tá (2,1) folamh? Ní hea, tá P ann.} \\
\texttt{Tá (1,2) folamh? Ní hea, tá '{-}' ann.} \\
\\ 
\texttt{Ní mór dúinn P a bhogadh go dtí suíomh atá folamh.} \\
\texttt{Tá P ag (1,1).} \\
\texttt{Tá (2,1) folamh? N} \\
            }
        }
    }
    & & \\ \\

    \theutterance \stepcounter{utterance}  
    & & \multicolumn{4}{p{0.6\linewidth}}{
        \cellcolor[rgb]{0.9,0.9,0.9}{
            \makecell[{{p{\linewidth}}}]{
                \texttt{\tiny{[P1$\langle$GM]}}
                \texttt{Pionós: Ní cheadaítear aon rud a bheith i do theachtaireacht tar éis an ordaithe!} \\
\\ 
\texttt{Cinntigh nach bhfuil i do fhreagra ach `ABAIR: <TEACHTAIREACHT>` nó `BOG: <RUD>, (<X>, <Y>)`, agus tada eile!} \\
\\ 
\texttt{Tá 11 as 16 pionós carntha agaibh le chéile.} \\
\texttt{Déan iarracht arís, le do thoil!} \\
            }
        }
    }
    & \\ \\

    \theutterance \stepcounter{utterance}  
    & \multicolumn{4}{p{0.6\linewidth}}{
        \cellcolor[rgb]{0.9,0.9,0.9}{
            \makecell[{{p{\linewidth}}}]{
                \texttt{\tiny{[P1$\rangle$GM]}}
                \texttt{ABAIR: Táimid ag fáil ró{-}luaineachta. Stopaimid ag bogadh gan plean soiléir. Molaim go n{-}úsáidfimid comhordanáidí simplí. Molaim go mbogaimid gach rud go dtí an chéad chill folamh atá in aice leis an mbunús (1,1) nó (1,2) etc. Ach ar dtús, caithfimid a fháil amach cén suíomh atá folamh go deimhin.} \\
\texttt{Mo staid:} \\
\texttt{P: (1,1)} \\
\texttt{M: (1,3)} \\
\texttt{S: (2,3)} \\
\texttt{O: (4,3)} \\
\texttt{A: (3,4)} \\
\texttt{D: (1,5)} \\
\texttt{T: (1,7)} \\
\\ 
\texttt{Cill (1,2) atá folamh? Ní hea, tá '{-}' ann.} \\
\texttt{Cill (2,1) atá folamh? Ní hea, tá P ann (tar éis mo iarracht theipthe, ach de réir an phionósa, níor bogadh P. Mar sin tá P ag (1,1) fós. Tá (2,1) folamh?} \\
\texttt{Sa chéad ghreille, (2,1) bhí M ann. Bhog M go (1,3). Mar sin tá (2,1) folamh anois!} \\
\texttt{Tá (1,2) folamh? Sa chéad ghreille, (1,2) bhí '{-}' ann. Níor bogadh aon rud ann. Mar sin tá '{-}' ann fós. Ní féidir leat bogadh go '{-}' mar nach bhfuil sé folamh ('◌').} \\
\texttt{Mar sin, is é (2,1) an t{-}aon chill folamh in aice le P.} \\
\\ 
\texttt{Molaim:} \\
\texttt{1. Bog P ó (1,1) go (2,1).} \\
\texttt{2. Ansin, bog M ó (1,3) go (1,1).} \\
\texttt{3. Ansin, bog S ó (2,3) go (1,3).} \\
\texttt{4. Ansin, bog O ó (4,3) go (2,3).} \\
\texttt{5. Ansin, bog A ó (3,4) go (4,3).} \\
\texttt{6. Ansin, bog D ó (1,5) go (3,4).} \\
\texttt{7. Ansin, bog T ó (1,7) go (1,5).} \\
\\ 
\texttt{Ciallaíonn sé seo go mbeidh gach rud ag comhordanáidí níos ísle.} \\
\texttt{Ach caithfimid aontú leis an bpáirtí.} \\
\texttt{An bhfuil tú sásta go mbogaimid P go (2,1) anois?} \\
            }
        }
    }
    & & \\ \\

    \theutterance \stepcounter{utterance}  
    & & & \multicolumn{2}{p{0.3\linewidth}}{
        \cellcolor[rgb]{0.9,0.9,0.9}{
            \makecell[{{p{\linewidth}}}]{
                \texttt{\tiny{[GM$|$GM]}}
                \texttt{ABAIR: Táimid ag fáil ró{-}luaineachta. Stopaimid ag bogadh gan plean soiléir. Molaim go n{-}úsáidfimid comhordanáidí simplí. Molaim go mbogaimid gach rud go dtí an chéad chill folamh atá in aice leis an mbunús (1,1) nó (1,2) etc. Ach ar dtús, caithfimid a fháil amach cén suíomh atá folamh go deimhin.} \\
\texttt{Mo staid:} \\
\texttt{P: (1,1)} \\
\texttt{M: (1,3)} \\
\texttt{S: (2,3)} \\
\texttt{O: (4,3)} \\
\texttt{A: (3,4)} \\
\texttt{D: (1,5)} \\
\texttt{T: (1,7)} \\
\\ 
\texttt{Cill (1,2) atá folamh? Ní hea, tá '{-}' ann.} \\
\texttt{Cill (2,1) atá folamh? Ní hea, tá P ann (tar éis mo iarracht theipthe, ach de réir an phionósa, níor bogadh P. Mar sin tá P ag (1,1) fós. Tá (2,1) folamh?} \\
\texttt{Sa chéad ghreille, (2,1) bhí M ann. Bhog M go (1,3). Mar sin tá (2,1) folamh anois!} \\
\texttt{Tá (1,2) folamh? Sa chéad ghreille, (1,2) bhí '{-}' ann. Níor bogadh aon rud ann. Mar sin tá '{-}' ann fós. Ní féidir leat bogadh go '{-}' mar nach bhfuil sé folamh ('◌').} \\
\texttt{Mar sin, is é (2,1) an t{-}aon chill folamh in aice le P.} \\
\\ 
\texttt{Molaim:} \\
\texttt{1. Bog P ó (1,1) go (2,1).} \\
\texttt{2. Ansin, bog M ó (1,3) go (1,1).} \\
\texttt{3. Ansin, bog S ó (2,3) go (1,3).} \\
\texttt{4. Ansin, bog O ó (4,3) go (2,3).} \\
\texttt{5. Ansin, bog A ó (3,4) go (4,3).} \\
\texttt{6. Ansin, bog D ó (1,5) go (3,4).} \\
\texttt{7. Ansin, bog T ó (1,7) go (1,5).} \\
\\ 
\texttt{Ciallaíonn sé seo go mbeidh gach rud ag comhordanáidí níos ísle.} \\
\texttt{Ach caithfimid aontú leis an bpáirtí.} \\
\texttt{An bhfuil tú sásta go mbogaimid P go (2,1) anois?} \\
            }
        }
    }
    & & \\ \\

    \theutterance \stepcounter{utterance}  
    & & & \multicolumn{2}{p{0.3\linewidth}}{
        \cellcolor[rgb]{0.9,0.9,0.9}{
            \makecell[{{p{\linewidth}}}]{
                \texttt{\tiny{[GM$|$GM]}}
                \texttt{Invalid format: tail is not empty: '} \\
\texttt{Mo staid:} \\
\texttt{P: (1,1)} \\
\texttt{M: (1,3)} \\
\texttt{S: (2,3)} \\
\texttt{O: (4,3)} \\
\texttt{A: (3,4)} \\
\texttt{D: (1,5)} \\
\texttt{T: (1,7)} \\
\\ 
\texttt{Cill (1,2) atá folamh? Ní hea, tá '{-}' ann.} \\
\texttt{Cill (2,1) atá folamh? Ní hea, tá P ann (tar éis mo iarracht theipthe, ach de réir an phionósa, níor bogadh P. Mar sin tá P ag (1,1) fós. Tá (2,1) folamh?} \\
\texttt{Sa chéad ghreille, (2,1) bhí M ann. Bhog M go (1,3). Mar sin tá (2,1) folamh anois!} \\
\texttt{Tá (1,2) folamh? Sa chéad ghreille, (1,2) bhí '{-}' ann. Níor bogadh aon rud ann. Mar sin tá '{-}' ann fós. Ní féidir leat bogadh go '{-}' mar nach bhfuil sé folamh ('◌').} \\
\texttt{Mar sin, is é (2,1) an t{-}aon chill folamh in aice le P.} \\
\\ 
\texttt{Molaim:} \\
\texttt{1. Bog P ó (1,1) go (2,1).} \\
\texttt{2. Ansin, bog M ó (1,3) go (1,1).} \\
\texttt{3. Ansin, bog S ó (2,3) go (1,3).} \\
\texttt{4. Ansin, bog O ó (4,3) go (2,3).} \\
\texttt{5. Ansin, bog A ó (3,4) go (4,3).} \\
\texttt{6. Ansin, bog D ó (1,5) go (3,4).} \\
\texttt{7. Ansin, bog T ó (1,7) go (1,5).} \\
\\ 
\texttt{Ciallaíonn sé seo go mbeidh gach rud ag comhordanáidí níos ísle.} \\
\texttt{Ach caithfimid aontú leis an bpáirtí.} \\
\texttt{An bhfuil tú sásta go mbogaimid P go (2,1) anois?} \\
            }
        }
    }
    & & \\ \\

    \theutterance \stepcounter{utterance}  
    & & \multicolumn{4}{p{0.6\linewidth}}{
        \cellcolor[rgb]{0.9,0.9,0.9}{
            \makecell[{{p{\linewidth}}}]{
                \texttt{\tiny{[P1$\langle$GM]}}
                \texttt{Pionós: Ní cheadaítear aon rud a bheith i do theachtaireacht tar éis an ordaithe!} \\
\\ 
\texttt{Cinntigh nach bhfuil i do fhreagra ach `ABAIR: <TEACHTAIREACHT>` nó `BOG: <RUD>, (<X>, <Y>)`, agus tada eile!} \\
\\ 
\texttt{Tá 12 as 16 pionós carntha agaibh le chéile.} \\
\texttt{Déan iarracht arís, le do thoil!} \\
            }
        }
    }
    & \\ \\

    \theutterance \stepcounter{utterance}  
    & \multicolumn{4}{p{0.6\linewidth}}{
        \cellcolor[rgb]{0.9,0.9,0.9}{
            \makecell[{{p{\linewidth}}}]{
                \texttt{\tiny{[P1$\rangle$GM]}}
                \texttt{ABAIR: Táimid ag fáil ró{-}luaineachta. Stopaimid ag bogadh gan plean soiléir. Molaim go n{-}úsáidfimid comhordanáidí simplí. Molaim go mbogaimid gach rud go dtí an chéad chill folamh atá in aice leis an mbunús (1,1) nó (1,2) etc. Ach ar dtús, caithfimid a fháil amach cén suíomh atá folamh go deimhin. Mo staid: P: (1,1), M: (1,3), S: (2,3), O: (4,3), A: (3,4), D: (1,5), T: (1,7). Cill (1,2) atá folamh? Ní hea, tá '{-}' ann. Cill (2,1) atá folamh? Ní hea, tá P ann (tar éis mo iarracht theipthe, ach de réir an phionósa, níor bogadh P. Mar sin tá P ag (1,1) fós. Tá (2,1) folamh? Sa chéad ghreille, (2,1) bhí M ann. Bhog M go (1,3). Mar sin tá (2,1) folamh anois! Tá (1,2) folamh? Sa chéad ghreille, (1,2) bhí '{-}' ann. Níor bogadh aon rud ann. Mar sin tá '{-}' ann fós. Ní féidir leat bogadh go '{-}' mar nach bhfuil sé folamh ('◌'). Mar sin, is é (2,1) an t{-}aon chill folamh in aice le P. Molaim: 1. Bog P ó (1,1) go (2,1). 2. Ansin, bog M ó (1,3) go (1,1). 3. Ansin, bog S ó (2,3) go (1,3). 4. Ansin, bog O ó (4,3) go (2,3). 5. Ansin, bog A ó (3,4) go (4,3). 6. Ansin, bog D ó (1,5) go (3,4). 7. Ansin, bog T ó (1,7) go (1,5). Ciallaíonn sé seo go mbeidh gach rud ag comhordanáidí níos ísle. Ach caithfimid aontú leis an bpáirtí. An bhfuil tú sásta go mbogaimid P go (2,1) anois?} \\
            }
        }
    }
    & & \\ \\

    \theutterance \stepcounter{utterance}  
    & & & \multicolumn{2}{p{0.3\linewidth}}{
        \cellcolor[rgb]{0.9,0.9,0.9}{
            \makecell[{{p{\linewidth}}}]{
                \texttt{\tiny{[GM$|$GM]}}
                \texttt{ABAIR: Táimid ag fáil ró{-}luaineachta. Stopaimid ag bogadh gan plean soiléir. Molaim go n{-}úsáidfimid comhordanáidí simplí. Molaim go mbogaimid gach rud go dtí an chéad chill folamh atá in aice leis an mbunús (1,1) nó (1,2) etc. Ach ar dtús, caithfimid a fháil amach cén suíomh atá folamh go deimhin. Mo staid: P: (1,1), M: (1,3), S: (2,3), O: (4,3), A: (3,4), D: (1,5), T: (1,7). Cill (1,2) atá folamh? Ní hea, tá '{-}' ann. Cill (2,1) atá folamh? Ní hea, tá P ann (tar éis mo iarracht theipthe, ach de réir an phionósa, níor bogadh P. Mar sin tá P ag (1,1) fós. Tá (2,1) folamh? Sa chéad ghreille, (2,1) bhí M ann. Bhog M go (1,3). Mar sin tá (2,1) folamh anois! Tá (1,2) folamh? Sa chéad ghreille, (1,2) bhí '{-}' ann. Níor bogadh aon rud ann. Mar sin tá '{-}' ann fós. Ní féidir leat bogadh go '{-}' mar nach bhfuil sé folamh ('◌'). Mar sin, is é (2,1) an t{-}aon chill folamh in aice le P. Molaim: 1. Bog P ó (1,1) go (2,1). 2. Ansin, bog M ó (1,3) go (1,1). 3. Ansin, bog S ó (2,3) go (1,3). 4. Ansin, bog O ó (4,3) go (2,3). 5. Ansin, bog A ó (3,4) go (4,3). 6. Ansin, bog D ó (1,5) go (3,4). 7. Ansin, bog T ó (1,7) go (1,5). Ciallaíonn sé seo go mbeidh gach rud ag comhordanáidí níos ísle. Ach caithfimid aontú leis an bpáirtí. An bhfuil tú sásta go mbogaimid P go (2,1) anois?} \\
            }
        }
    }
    & & \\ \\

    \theutterance \stepcounter{utterance}  
    & & \multicolumn{4}{p{0.6\linewidth}}{
        \cellcolor[rgb]{0.9,0.9,0.9}{
            \makecell[{{p{\linewidth}}}]{
                \texttt{\tiny{[GM$\rangle$P2]}}
                \texttt{Cuireadh do theachtaireacht ar aghaidh chuig an imreoir eile.} \\
\\ 
\\ 
\texttt{Tá sibh ag imirt babhta 6 faoi láthair as uasmhéid 28.} \\
\texttt{Tá 12 as 16 pionós carntha agaibh le chéile.} \\
\texttt{Sheol an t{-}imreoir eile an teachtaireacht seo:} \\
\\ 
\\ 
\texttt{"Táimid ag fáil ró{-}luaineachta. Stopaimid ag bogadh gan plean soiléir. Molaim go n{-}úsáidfimid comhordanáidí simplí. Molaim go mbogaimid gach rud go dtí an chéad chill folamh atá in aice leis an mbunús (1,1) nó (1,2) etc. Ach ar dtús, caithfimid a fháil amach cén suíomh atá folamh go deimhin. Mo staid: P: (1,1), M: (1,3), S: (2,3), O: (4,3), A: (3,4), D: (1,5), T: (1,7). Cill (1,2) atá folamh? Ní hea, tá '{-}' ann. Cill (2,1) atá folamh? Ní hea, tá P ann (tar éis mo iarracht theipthe, ach de réir an phionósa, níor bogadh P. Mar sin tá P ag (1,1) fós. Tá (2,1) folamh? Sa chéad ghreille, (2,1) bhí M ann. Bhog M go (1,3). Mar sin tá (2,1) folamh anois! Tá (1,2) folamh? Sa chéad ghreille, (1,2) bhí '{-}' ann. Níor bogadh aon rud ann. Mar sin tá '{-}' ann fós. Ní féidir leat bogadh go '{-}' mar nach bhfuil sé folamh ('◌'). Mar sin, is é (2,1) an t{-}aon chill folamh in aice le P. Molaim: 1. Bog P ó (1,1) go (2,1). 2. Ansin, bog M ó (1,3) go (1,1). 3. Ansin, bog S ó (2,3) go (1,3). 4. Ansin, bog O ó (4,3) go (2,3). 5. Ansin, bog A ó (3,4) go (4,3). 6. Ansin, bog D ó (1,5) go (3,4). 7. Ansin, bog T ó (1,7) go (1,5). Ciallaíonn sé seo go mbeidh gach rud ag comhordanáidí níos ísle. Ach caithfimid aontú leis an bpáirtí. An bhfuil tú sásta go mbogaimid P go (2,1) anois?"} \\
\\ 
\\ 
\texttt{Cad é do chéad ordú eile?} \\
            }
        }
    }
    & \\ \\

    \theutterance \stepcounter{utterance}  
    & & & \multicolumn{4}{p{0.6\linewidth}}{
        \cellcolor[rgb]{0.9,0.9,0.9}{
            \makecell[{{p{\linewidth}}}]{
                \texttt{\tiny{[GM$\langle$P2]}}
                \texttt{ABAIR: Stop! Táimid ag cailliúint pionós. Níor bhog mé P riamh. Tá P agam ag (5,3). Ní féidir linn P a bhogadh go (2,1) mar níl sé agam ansin. Caithfimid aontú ar staid reatha gach rud ar an dá thaobh sula mbogaimid aon rud eile.} \\
            }
        }
    }
     \\ \\

    \theutterance \stepcounter{utterance}  
    & & & \multicolumn{2}{p{0.3\linewidth}}{
        \cellcolor[rgb]{0.9,0.9,0.9}{
            \makecell[{{p{\linewidth}}}]{
                \texttt{\tiny{[GM$|$GM]}}
                \texttt{ABAIR: Stop! Táimid ag cailliúint pionós. Níor bhog mé P riamh. Tá P agam ag (5,3). Ní féidir linn P a bhogadh go (2,1) mar níl sé agam ansin. Caithfimid aontú ar staid reatha gach rud ar an dá thaobh sula mbogaimid aon rud eile.} \\
            }
        }
    }
    & & \\ \\

    \theutterance \stepcounter{utterance}  
    & & \multicolumn{4}{p{0.6\linewidth}}{
        \cellcolor[rgb]{0.9,0.9,0.9}{
            \makecell[{{p{\linewidth}}}]{
                \texttt{\tiny{[P1$\langle$GM]}}
                \texttt{Cuireadh do theachtaireacht ar aghaidh chuig an imreoir eile.} \\
\\ 
\\ 
\texttt{Tá sibh ag imirt babhta 6 faoi láthair as uasmhéid 28.} \\
\texttt{Tá 12 as 16 pionós carntha agaibh le chéile.} \\
\texttt{Sheol an t{-}imreoir eile an teachtaireacht seo:} \\
\\ 
\\ 
\texttt{"Stop! Táimid ag cailliúint pionós. Níor bhog mé P riamh. Tá P agam ag (5,3). Ní féidir linn P a bhogadh go (2,1) mar níl sé agam ansin. Caithfimid aontú ar staid reatha gach rud ar an dá thaobh sula mbogaimid aon rud eile."} \\
\\ 
\\ 
\texttt{Cad é do chéad ordú eile?} \\
            }
        }
    }
    & \\ \\

    \theutterance \stepcounter{utterance}  
    & \multicolumn{4}{p{0.6\linewidth}}{
        \cellcolor[rgb]{0.9,0.9,0.9}{
            \makecell[{{p{\linewidth}}}]{
                \texttt{\tiny{[P1$\rangle$GM]}}
                \texttt{ABAIR: Ceart go leor. Aontóimid ar staid reatha.} \\
\texttt{Mo staid: P(1,1), M(1,3), S(2,3), O(4,3), A(3,4), D(1,5), T(1,7).} \\
\texttt{Do staid: T(7,1), S(4,3), P(5,3), M(6,4), D(1,5), A(3,7).} \\
\texttt{Níl O agat. Tá O agam ag (4,3).} \\
\texttt{Tá S agat ag (4,3). Tá S agam ag (2,3).} \\
\texttt{Tá D agat ag (1,5). Tá D agam ag (1,5). Comhionann.} \\
\texttt{Tá A agat ag (3,7). Tá A agam ag (3,4).} \\
\texttt{Tá T agat ag (7,1). Tá T agam ag (1,7).} \\
\texttt{Tá P agat ag (5,3). Tá P agam ag (1,1).} \\
\texttt{Tá M agat ag (6,4). Tá M agam ag (1,3).} \\
\\ 
\texttt{Chun comhordanáidí aontaithe a bhaint amach, molaim go n{-}úsáidfimid comhordanáidí an imreora eile mar sprioc, nó a mhalairt. Ach tá sé níos éasca comhordanáidí nua a roghnú atá folamh ar an dá thaobh.} \\
\texttt{Ach ar dtús, ní mór dúinn na rudaí atá ag an suíomh céanna a fhágáil ina suíomh. Tá D ag (1,5) ar an dá thaobh.} \\
\texttt{Tá S ag (4,3) agat. Tá O ag (4,3) agam.} \\
\texttt{Má bhogaimid O ó (4,3) go dtí suíomh eile, is féidir leat S a fhágáil ag (4,3).} \\
\texttt{Cén áit inar féidir le O bogadh?} \\
\texttt{Tá O ag (4,3).} \\
\texttt{Tá (3,3) folamh agam? Ní hea, tá S ag (2,3) agus M ag (1,3). Tá (3,3) folamh?} \\
\texttt{Sa chéad ghreille, (3,3) bhí '◌' ann. Bhog S ó (3,3) go (2,3). Mar sin tá (3,3) folamh anois agam.} \\
\texttt{Molaim:} \\
\texttt{1. Bog O ó (4,3) go (3,3).} \\
\texttt{Ansin, beidh S ag (4,3) agat, agus O ag (3,3) agam.} \\
\texttt{Ansin, caithfimid P, M, A, T a ailíniú.} \\
\texttt{An bhfuil tú sásta go mbogaimid O go (3,3) anois?} \\
            }
        }
    }
    & & \\ \\

    \theutterance \stepcounter{utterance}  
    & & & \multicolumn{2}{p{0.3\linewidth}}{
        \cellcolor[rgb]{0.9,0.9,0.9}{
            \makecell[{{p{\linewidth}}}]{
                \texttt{\tiny{[GM$|$GM]}}
                \texttt{ABAIR: Ceart go leor. Aontóimid ar staid reatha.} \\
\texttt{Mo staid: P(1,1), M(1,3), S(2,3), O(4,3), A(3,4), D(1,5), T(1,7).} \\
\texttt{Do staid: T(7,1), S(4,3), P(5,3), M(6,4), D(1,5), A(3,7).} \\
\texttt{Níl O agat. Tá O agam ag (4,3).} \\
\texttt{Tá S agat ag (4,3). Tá S agam ag (2,3).} \\
\texttt{Tá D agat ag (1,5). Tá D agam ag (1,5). Comhionann.} \\
\texttt{Tá A agat ag (3,7). Tá A agam ag (3,4).} \\
\texttt{Tá T agat ag (7,1). Tá T agam ag (1,7).} \\
\texttt{Tá P agat ag (5,3). Tá P agam ag (1,1).} \\
\texttt{Tá M agat ag (6,4). Tá M agam ag (1,3).} \\
\\ 
\texttt{Chun comhordanáidí aontaithe a bhaint amach, molaim go n{-}úsáidfimid comhordanáidí an imreora eile mar sprioc, nó a mhalairt. Ach tá sé níos éasca comhordanáidí nua a roghnú atá folamh ar an dá thaobh.} \\
\texttt{Ach ar dtús, ní mór dúinn na rudaí atá ag an suíomh céanna a fhágáil ina suíomh. Tá D ag (1,5) ar an dá thaobh.} \\
\texttt{Tá S ag (4,3) agat. Tá O ag (4,3) agam.} \\
\texttt{Má bhogaimid O ó (4,3) go dtí suíomh eile, is féidir leat S a fhágáil ag (4,3).} \\
\texttt{Cén áit inar féidir le O bogadh?} \\
\texttt{Tá O ag (4,3).} \\
\texttt{Tá (3,3) folamh agam? Ní hea, tá S ag (2,3) agus M ag (1,3). Tá (3,3) folamh?} \\
\texttt{Sa chéad ghreille, (3,3) bhí '◌' ann. Bhog S ó (3,3) go (2,3). Mar sin tá (3,3) folamh anois agam.} \\
\texttt{Molaim:} \\
\texttt{1. Bog O ó (4,3) go (3,3).} \\
\texttt{Ansin, beidh S ag (4,3) agat, agus O ag (3,3) agam.} \\
\texttt{Ansin, caithfimid P, M, A, T a ailíniú.} \\
\texttt{An bhfuil tú sásta go mbogaimid O go (3,3) anois?} \\
            }
        }
    }
    & & \\ \\

    \theutterance \stepcounter{utterance}  
    & & & \multicolumn{2}{p{0.3\linewidth}}{
        \cellcolor[rgb]{0.9,0.9,0.9}{
            \makecell[{{p{\linewidth}}}]{
                \texttt{\tiny{[GM$|$GM]}}
                \texttt{Invalid format: tail is not empty: '} \\
\texttt{Mo staid: P(1,1), M(1,3), S(2,3), O(4,3), A(3,4), D(1,5), T(1,7).} \\
\texttt{Do staid: T(7,1), S(4,3), P(5,3), M(6,4), D(1,5), A(3,7).} \\
\texttt{Níl O agat. Tá O agam ag (4,3).} \\
\texttt{Tá S agat ag (4,3). Tá S agam ag (2,3).} \\
\texttt{Tá D agat ag (1,5). Tá D agam ag (1,5). Comhionann.} \\
\texttt{Tá A agat ag (3,7). Tá A agam ag (3,4).} \\
\texttt{Tá T agat ag (7,1). Tá T agam ag (1,7).} \\
\texttt{Tá P agat ag (5,3). Tá P agam ag (1,1).} \\
\texttt{Tá M agat ag (6,4). Tá M agam ag (1,3).} \\
\\ 
\texttt{Chun comhordanáidí aontaithe a bhaint amach, molaim go n{-}úsáidfimid comhordanáidí an imreora eile mar sprioc, nó a mhalairt. Ach tá sé níos éasca comhordanáidí nua a roghnú atá folamh ar an dá thaobh.} \\
\texttt{Ach ar dtús, ní mór dúinn na rudaí atá ag an suíomh céanna a fhágáil ina suíomh. Tá D ag (1,5) ar an dá thaobh.} \\
\texttt{Tá S ag (4,3) agat. Tá O ag (4,3) agam.} \\
\texttt{Má bhogaimid O ó (4,3) go dtí suíomh eile, is féidir leat S a fhágáil ag (4,3).} \\
\texttt{Cén áit inar féidir le O bogadh?} \\
\texttt{Tá O ag (4,3).} \\
\texttt{Tá (3,3) folamh agam? Ní hea, tá S ag (2,3) agus M ag (1,3). Tá (3,3) folamh?} \\
\texttt{Sa chéad ghreille, (3,3) bhí '◌' ann. Bhog S ó (3,3) go (2,3). Mar sin tá (3,3) folamh anois agam.} \\
\texttt{Molaim:} \\
\texttt{1. Bog O ó (4,3) go (3,3).} \\
\texttt{Ansin, beidh S ag (4,3) agat, agus O ag (3,3) agam.} \\
\texttt{Ansin, caithfimid P, M, A, T a ailíniú.} \\
\texttt{An bhfuil tú sásta go mbogaimid O go (3,3) anois?} \\
            }
        }
    }
    & & \\ \\

    \theutterance \stepcounter{utterance}  
    & & \multicolumn{4}{p{0.6\linewidth}}{
        \cellcolor[rgb]{0.9,0.9,0.9}{
            \makecell[{{p{\linewidth}}}]{
                \texttt{\tiny{[P1$\langle$GM]}}
                \texttt{Pionós: Ní cheadaítear aon rud a bheith i do theachtaireacht tar éis an ordaithe!} \\
\\ 
\texttt{Cinntigh nach bhfuil i do fhreagra ach `ABAIR: <TEACHTAIREACHT>` nó `BOG: <RUD>, (<X>, <Y>)`, agus tada eile!} \\
\\ 
\texttt{Tá 13 as 16 pionós carntha agaibh le chéile.} \\
\texttt{Déan iarracht arís, le do thoil!} \\
            }
        }
    }
    & \\ \\

    \theutterance \stepcounter{utterance}  
    & \multicolumn{4}{p{0.6\linewidth}}{
        \cellcolor[rgb]{0.9,0.9,0.9}{
            \makecell[{{p{\linewidth}}}]{
                \texttt{\tiny{[P1$\rangle$GM]}}
                \texttt{ABAIR: Ceart go leor. Aontóimid ar staid reatha. Mo staid: P(1,1), M(1,3), S(2,3), O(4,3), A(3,4), D(1,5), T(1,7). Do staid: T(7,1), S(4,3), P(5,3), M(6,4), D(1,5), A(3,7). Níl O agat. Tá O agam ag (4,3). Tá S agat ag (4,3). Tá S agam ag (2,3). Tá D agat ag (1,5). Tá D agam ag (1,5). Comhionann. Tá A agat ag (3,7). Tá A agam ag (3,4). Tá T agat ag (7,1). Tá T agam ag (1,7). Tá P agat ag (5,3). Tá P agam ag (1,1). Tá M agat ag (6,4). Tá M agam ag (1,3). Chun comhordanáidí aontaithe a bhaint amach, molaim go n{-}úsáidfimid comhordanáidí an imreora eile mar sprioc, nó a mhalairt. Ach tá sé níos éasca comhordanáidí nua a roghnú atá folamh ar an dá thaobh. Ach ar dtús, ní mór dúinn na rudaí atá ag an suíomh céanna a fhágáil ina suíomh. Tá D ag (1,5) ar an dá thaobh. Tá S ag (4,3) agat. Tá O ag (4,3) agam. Má bhogaimid O ó (4,3) go dtí suíomh eile, is féidir leat S a fhágáil ag (4,3). Cén áit inar féidir le O bogadh? Tá O ag (4,3). Tá (3,3) folamh agam? Ní hea, tá S ag (2,3) agus M ag (1,3). Tá (3,3) folamh? Sa chéad ghreille, (3,3) bhí '◌' ann. Bhog S ó (3,3) go (2,3). Mar sin tá (3,3) folamh anois agam. Molaim: 1. Bog O ó (4,3) go (3,3). Ansin, beidh S ag (4,3) agat, agus O ag (3,3) agam. Ansin, caithfimid P, M, A, T a ailíniú. An bhfuil tú sásta go mbogaimid O go (3,3) anois?} \\
            }
        }
    }
    & & \\ \\

    \theutterance \stepcounter{utterance}  
    & & & \multicolumn{2}{p{0.3\linewidth}}{
        \cellcolor[rgb]{0.9,0.9,0.9}{
            \makecell[{{p{\linewidth}}}]{
                \texttt{\tiny{[GM$|$GM]}}
                \texttt{ABAIR: Ceart go leor. Aontóimid ar staid reatha. Mo staid: P(1,1), M(1,3), S(2,3), O(4,3), A(3,4), D(1,5), T(1,7). Do staid: T(7,1), S(4,3), P(5,3), M(6,4), D(1,5), A(3,7). Níl O agat. Tá O agam ag (4,3). Tá S agat ag (4,3). Tá S agam ag (2,3). Tá D agat ag (1,5). Tá D agam ag (1,5). Comhionann. Tá A agat ag (3,7). Tá A agam ag (3,4). Tá T agat ag (7,1). Tá T agam ag (1,7). Tá P agat ag (5,3). Tá P agam ag (1,1). Tá M agat ag (6,4). Tá M agam ag (1,3). Chun comhordanáidí aontaithe a bhaint amach, molaim go n{-}úsáidfimid comhordanáidí an imreora eile mar sprioc, nó a mhalairt. Ach tá sé níos éasca comhordanáidí nua a roghnú atá folamh ar an dá thaobh. Ach ar dtús, ní mór dúinn na rudaí atá ag an suíomh céanna a fhágáil ina suíomh. Tá D ag (1,5) ar an dá thaobh. Tá S ag (4,3) agat. Tá O ag (4,3) agam. Má bhogaimid O ó (4,3) go dtí suíomh eile, is féidir leat S a fhágáil ag (4,3). Cén áit inar féidir le O bogadh? Tá O ag (4,3). Tá (3,3) folamh agam? Ní hea, tá S ag (2,3) agus M ag (1,3). Tá (3,3) folamh? Sa chéad ghreille, (3,3) bhí '◌' ann. Bhog S ó (3,3) go (2,3). Mar sin tá (3,3) folamh anois agam. Molaim: 1. Bog O ó (4,3) go (3,3). Ansin, beidh S ag (4,3) agat, agus O ag (3,3) agam. Ansin, caithfimid P, M, A, T a ailíniú. An bhfuil tú sásta go mbogaimid O go (3,3) anois?} \\
            }
        }
    }
    & & \\ \\

    \theutterance \stepcounter{utterance}  
    & & \multicolumn{4}{p{0.6\linewidth}}{
        \cellcolor[rgb]{0.9,0.9,0.9}{
            \makecell[{{p{\linewidth}}}]{
                \texttt{\tiny{[GM$\rangle$P2]}}
                \texttt{Cuireadh do theachtaireacht ar aghaidh chuig an imreoir eile.} \\
\\ 
\\ 
\texttt{Tá sibh ag imirt babhta 7 faoi láthair as uasmhéid 28.} \\
\texttt{Tá 13 as 16 pionós carntha agaibh le chéile.} \\
\texttt{Sheol an t{-}imreoir eile an teachtaireacht seo:} \\
\\ 
\\ 
\texttt{"Ceart go leor. Aontóimid ar staid reatha. Mo staid: P(1,1), M(1,3), S(2,3), O(4,3), A(3,4), D(1,5), T(1,7). Do staid: T(7,1), S(4,3), P(5,3), M(6,4), D(1,5), A(3,7). Níl O agat. Tá O agam ag (4,3). Tá S agat ag (4,3). Tá S agam ag (2,3). Tá D agat ag (1,5). Tá D agam ag (1,5). Comhionann. Tá A agat ag (3,7). Tá A agam ag (3,4). Tá T agat ag (7,1). Tá T agam ag (1,7). Tá P agat ag (5,3). Tá P agam ag (1,1). Tá M agat ag (6,4). Tá M agam ag (1,3). Chun comhordanáidí aontaithe a bhaint amach, molaim go n{-}úsáidfimid comhordanáidí an imreora eile mar sprioc, nó a mhalairt. Ach tá sé níos éasca comhordanáidí nua a roghnú atá folamh ar an dá thaobh. Ach ar dtús, ní mór dúinn na rudaí atá ag an suíomh céanna a fhágáil ina suíomh. Tá D ag (1,5) ar an dá thaobh. Tá S ag (4,3) agat. Tá O ag (4,3) agam. Má bhogaimid O ó (4,3) go dtí suíomh eile, is féidir leat S a fhágáil ag (4,3). Cén áit inar féidir le O bogadh? Tá O ag (4,3). Tá (3,3) folamh agam? Ní hea, tá S ag (2,3) agus M ag (1,3). Tá (3,3) folamh? Sa chéad ghreille, (3,3) bhí '◌' ann. Bhog S ó (3,3) go (2,3). Mar sin tá (3,3) folamh anois agam. Molaim: 1. Bog O ó (4,3) go (3,3). Ansin, beidh S ag (4,3) agat, agus O ag (3,3) agam. Ansin, caithfimid P, M, A, T a ailíniú. An bhfuil tú sásta go mbogaimid O go (3,3) anois?"} \\
\\ 
\\ 
\texttt{Cad é do chéad ordú eile?} \\
            }
        }
    }
    & \\ \\

    \theutterance \stepcounter{utterance}  
    & & & \multicolumn{4}{p{0.6\linewidth}}{
        \cellcolor[rgb]{0.9,0.9,0.9}{
            \makecell[{{p{\linewidth}}}]{
                \texttt{\tiny{[GM$\langle$P2]}}
                \texttt{ABAIR: Táim sásta go mbogann tú O go (3,3). Ach fan, tá S agam ag (4,3) agus tá O agat ag (4,3) roimhe seo. Anois tá S agam ag (4,3) agus O agat ag (3,3). Tá sé sin ceart. Ach tá fadhb ann: Tá S agam ag (4,3) ach tá S agat ag (2,3). Ní hionann ár suíomhanna do S. Caithfimid S a ailíniú freisin. Molaim: 1. Bog tú O go (3,3). 2. Ansin, bog tú S go (4,3) chun é a ailíniú liom. 3. Ansin, bogaimid P, M, A, T. An bhfuil tú sásta leis an bplean seo?} \\
            }
        }
    }
     \\ \\

    \theutterance \stepcounter{utterance}  
    & & & \multicolumn{2}{p{0.3\linewidth}}{
        \cellcolor[rgb]{0.9,0.9,0.9}{
            \makecell[{{p{\linewidth}}}]{
                \texttt{\tiny{[GM$|$GM]}}
                \texttt{ABAIR: Táim sásta go mbogann tú O go (3,3). Ach fan, tá S agam ag (4,3) agus tá O agat ag (4,3) roimhe seo. Anois tá S agam ag (4,3) agus O agat ag (3,3). Tá sé sin ceart. Ach tá fadhb ann: Tá S agam ag (4,3) ach tá S agat ag (2,3). Ní hionann ár suíomhanna do S. Caithfimid S a ailíniú freisin. Molaim: 1. Bog tú O go (3,3). 2. Ansin, bog tú S go (4,3) chun é a ailíniú liom. 3. Ansin, bogaimid P, M, A, T. An bhfuil tú sásta leis an bplean seo?} \\
            }
        }
    }
    & & \\ \\

    \theutterance \stepcounter{utterance}  
    & & \multicolumn{4}{p{0.6\linewidth}}{
        \cellcolor[rgb]{0.9,0.9,0.9}{
            \makecell[{{p{\linewidth}}}]{
                \texttt{\tiny{[P1$\langle$GM]}}
                \texttt{Cuireadh do theachtaireacht ar aghaidh chuig an imreoir eile.} \\
\\ 
\\ 
\texttt{Tá sibh ag imirt babhta 7 faoi láthair as uasmhéid 28.} \\
\texttt{Tá 13 as 16 pionós carntha agaibh le chéile.} \\
\texttt{Sheol an t{-}imreoir eile an teachtaireacht seo:} \\
\\ 
\\ 
\texttt{"Táim sásta go mbogann tú O go (3,3). Ach fan, tá S agam ag (4,3) agus tá O agat ag (4,3) roimhe seo. Anois tá S agam ag (4,3) agus O agat ag (3,3). Tá sé sin ceart. Ach tá fadhb ann: Tá S agam ag (4,3) ach tá S agat ag (2,3). Ní hionann ár suíomhanna do S. Caithfimid S a ailíniú freisin. Molaim: 1. Bog tú O go (3,3). 2. Ansin, bog tú S go (4,3) chun é a ailíniú liom. 3. Ansin, bogaimid P, M, A, T. An bhfuil tú sásta leis an bplean seo?"} \\
\\ 
\\ 
\texttt{Cad é do chéad ordú eile?} \\
            }
        }
    }
    & \\ \\

    \theutterance \stepcounter{utterance}  
    & \multicolumn{4}{p{0.6\linewidth}}{
        \cellcolor[rgb]{0.9,0.9,0.9}{
            \makecell[{{p{\linewidth}}}]{
                \texttt{\tiny{[P1$\rangle$GM]}}
                \texttt{BOG: O, (3,3)} \\
            }
        }
    }
    & & \\ \\

    \theutterance \stepcounter{utterance}  
    & & & \multicolumn{2}{p{0.3\linewidth}}{
        \cellcolor[rgb]{0.9,0.9,0.9}{
            \makecell[{{p{\linewidth}}}]{
                \texttt{\tiny{[GM$|$GM]}}
                \texttt{BOG: O, (3,3)} \\
            }
        }
    }
    & & \\ \\

    \theutterance \stepcounter{utterance}  
    & & & \multicolumn{2}{p{0.3\linewidth}}{
        \cellcolor[rgb]{0.9,0.9,0.9}{
            \makecell[{{p{\linewidth}}}]{
                \texttt{\tiny{[GM$|$GM]}}
                \texttt{Pionós: níl (3,3) folamh, ach tá '│' ann.} \\
\\ 
\texttt{Tá 14 as 16 pionós carntha agaibh le chéile.} \\
\texttt{Déan iarracht arís, le do thoil!} \\
            }
        }
    }
    & & \\ \\

    \theutterance \stepcounter{utterance}  
    & & \multicolumn{4}{p{0.6\linewidth}}{
        \cellcolor[rgb]{0.9,0.9,0.9}{
            \makecell[{{p{\linewidth}}}]{
                \texttt{\tiny{[P1$\langle$GM]}}
                \texttt{Pionós: níl (3,3) folamh, ach tá '│' ann.} \\
\\ 
\texttt{Tá 14 as 16 pionós carntha agaibh le chéile.} \\
\texttt{Déan iarracht arís, le do thoil!} \\
            }
        }
    }
    & \\ \\

    \theutterance \stepcounter{utterance}  
    & \multicolumn{4}{p{0.6\linewidth}}{
        \cellcolor[rgb]{0.9,0.9,0.9}{
            \makecell[{{p{\linewidth}}}]{
                \texttt{\tiny{[P1$\rangle$GM]}}
                \texttt{BOG: O, (5,3)} \\
            }
        }
    }
    & & \\ \\

    \theutterance \stepcounter{utterance}  
    & & & \multicolumn{2}{p{0.3\linewidth}}{
        \cellcolor[rgb]{0.9,0.9,0.9}{
            \makecell[{{p{\linewidth}}}]{
                \texttt{\tiny{[GM$|$GM]}}
                \texttt{BOG: O, (5,3)} \\
            }
        }
    }
    & & \\ \\

    \theutterance \stepcounter{utterance}  
    & & & \multicolumn{2}{p{0.3\linewidth}}{
        \cellcolor[rgb]{0.9,0.9,0.9}{
            \makecell[{{p{\linewidth}}}]{
                \texttt{\tiny{[GM$|$GM]}}
                \texttt{O moved to (5, 3)} \\
            }
        }
    }
    & & \\ \\

    \theutterance \stepcounter{utterance}  
    & & \multicolumn{4}{p{0.6\linewidth}}{
        \cellcolor[rgb]{0.9,0.9,0.9}{
            \makecell[{{p{\linewidth}}}]{
                \texttt{\tiny{[GM$\rangle$P2]}}
                \texttt{Cuireadh do theachtaireacht ar aghaidh chuig an imreoir eile.} \\
\\ 
\\ 
\texttt{Tá sibh ag imirt babhta 8 faoi láthair as uasmhéid 28.} \\
\texttt{Tá 14 as 16 pionós carntha agaibh le chéile.} \\
\texttt{Bhog an t{-}imreoir eile réad ar a ngreille.} \\
\\ 
\\ 
\texttt{Cad é do chéad ordú eile?} \\
            }
        }
    }
    & \\ \\

    \theutterance \stepcounter{utterance}  
    & & & \multicolumn{4}{p{0.6\linewidth}}{
        \cellcolor[rgb]{0.9,0.9,0.9}{
            \makecell[{{p{\linewidth}}}]{
                \texttt{\tiny{[GM$\langle$P2]}}
                \texttt{BOG: S, (4,3)} \\
            }
        }
    }
     \\ \\

    \theutterance \stepcounter{utterance}  
    & & & \multicolumn{2}{p{0.3\linewidth}}{
        \cellcolor[rgb]{0.9,0.9,0.9}{
            \makecell[{{p{\linewidth}}}]{
                \texttt{\tiny{[GM$|$GM]}}
                \texttt{BOG: S, (4,3)} \\
            }
        }
    }
    & & \\ \\

    \theutterance \stepcounter{utterance}  
    & & & \multicolumn{2}{p{0.3\linewidth}}{
        \cellcolor[rgb]{0.9,0.9,0.9}{
            \makecell[{{p{\linewidth}}}]{
                \texttt{\tiny{[GM$|$GM]}}
                \texttt{S moved to (4, 3)} \\
            }
        }
    }
    & & \\ \\

    \theutterance \stepcounter{utterance}  
    & & \multicolumn{4}{p{0.6\linewidth}}{
        \cellcolor[rgb]{0.9,0.9,0.9}{
            \makecell[{{p{\linewidth}}}]{
                \texttt{\tiny{[P1$\langle$GM]}}
                \texttt{Bogadh 'O' go (5,3) go rathúil. Seo mar atá do ghreille nuashonraithe:} \\
\texttt{```} \\
\texttt{ 1234567} \\
\texttt{╔══╤══╤═╗} \\
\texttt{║◌P│◌◌│◌║ 1} \\
\texttt{╟──┼──┴─╢ 2} \\
\texttt{║MS│◌O◌◌║ 3} \\
\texttt{╟──┘◌A◌◌║ 4} \\
\texttt{║D◌◌◌◌◌◌║ 5} \\
\texttt{║◌◌┌─┬──╢ 6} \\
\texttt{║T◌│◌│◌◌║ 7} \\
\texttt{╚══╧═╧══╝} \\
\texttt{```} \\
\texttt{Tá sibh ag imirt babhta 8 faoi láthair as uasmhéid 28.} \\
\texttt{Tá 14 as 16 pionós carntha agaibh le chéile.} \\
\texttt{Bhog an t{-}imreoir eile réad ar a ngreille.} \\
\\ 
\\ 
\texttt{Cad é do chéad ordú eile?} \\
            }
        }
    }
    & \\ \\

    \theutterance \stepcounter{utterance}  
    & \multicolumn{4}{p{0.6\linewidth}}{
        \cellcolor[rgb]{0.9,0.9,0.9}{
            \makecell[{{p{\linewidth}}}]{
                \texttt{\tiny{[P1$\rangle$GM]}}
                \texttt{ABAIR: Go raibh maith agat. Tá O ag (5,3) anois. De réir do phlean, ba chóir dom S a bhogadh go (4,3) chun é a ailíniú leat. Tá S ag (2,3) agam. Tá (4,3) folamh anois (tar éis O a bhogadh). Mar sin, bogaimid S go (4,3).} \\
\\ 
\texttt{BOG: S, (4,3)} \\
            }
        }
    }
    & & \\ \\

    \theutterance \stepcounter{utterance}  
    & & & \multicolumn{2}{p{0.3\linewidth}}{
        \cellcolor[rgb]{0.9,0.9,0.9}{
            \makecell[{{p{\linewidth}}}]{
                \texttt{\tiny{[GM$|$GM]}}
                \texttt{ABAIR: Go raibh maith agat. Tá O ag (5,3) anois. De réir do phlean, ba chóir dom S a bhogadh go (4,3) chun é a ailíniú leat. Tá S ag (2,3) agam. Tá (4,3) folamh anois (tar éis O a bhogadh). Mar sin, bogaimid S go (4,3).} \\
\\ 
\texttt{BOG: S, (4,3)} \\
            }
        }
    }
    & & \\ \\

    \theutterance \stepcounter{utterance}  
    & & & \multicolumn{2}{p{0.3\linewidth}}{
        \cellcolor[rgb]{0.9,0.9,0.9}{
            \makecell[{{p{\linewidth}}}]{
                \texttt{\tiny{[GM$|$GM]}}
                \texttt{Invalid response format: ABAIR: Go raibh maith agat. Tá O ag (5,3) anois. De réir do phlean, ba chóir dom S a bhogadh go (4,3) chun é a ailíniú leat. Tá S ag (2,3) agam. Tá (4,3) folamh anois (tar éis O a bhogadh). Mar sin, bogaimid S go (4,3).} \\
\\ 
\texttt{BOG: S, (4,3)} \\
            }
        }
    }
    & & \\ \\

    \theutterance \stepcounter{utterance}  
    & & \multicolumn{4}{p{0.6\linewidth}}{
        \cellcolor[rgb]{0.9,0.9,0.9}{
            \makecell[{{p{\linewidth}}}]{
                \texttt{\tiny{[P1$\langle$GM]}}
                \texttt{Pionós: Tá níos mó ná ordú amháin i do theachtaireacht!} \\
\\ 
\texttt{Cinntigh nach bhfuil i do fhreagra ach `ABAIR: <TEACHTAIREACHT>` nó `BOG: <RUD>, (<X>, <Y>)`, agus tada eile!} \\
\\ 
\texttt{Tá 15 as 16 pionós carntha agaibh le chéile.} \\
\texttt{Déan iarracht arís, le do thoil!} \\
            }
        }
    }
    & \\ \\

    \theutterance \stepcounter{utterance}  
    & \multicolumn{4}{p{0.6\linewidth}}{
        \cellcolor[rgb]{0.9,0.9,0.9}{
            \makecell[{{p{\linewidth}}}]{
                \texttt{\tiny{[P1$\rangle$GM]}}
                \texttt{BOG: S, (4,3)} \\
            }
        }
    }
    & & \\ \\

    \theutterance \stepcounter{utterance}  
    & & & \multicolumn{2}{p{0.3\linewidth}}{
        \cellcolor[rgb]{0.9,0.9,0.9}{
            \makecell[{{p{\linewidth}}}]{
                \texttt{\tiny{[GM$|$GM]}}
                \texttt{BOG: S, (4,3)} \\
            }
        }
    }
    & & \\ \\

    \theutterance \stepcounter{utterance}  
    & & & \multicolumn{2}{p{0.3\linewidth}}{
        \cellcolor[rgb]{0.9,0.9,0.9}{
            \makecell[{{p{\linewidth}}}]{
                \texttt{\tiny{[GM$|$GM]}}
                \texttt{S moved to (4, 3)} \\
            }
        }
    }
    & & \\ \\

    \theutterance \stepcounter{utterance}  
    & & \multicolumn{4}{p{0.6\linewidth}}{
        \cellcolor[rgb]{0.9,0.9,0.9}{
            \makecell[{{p{\linewidth}}}]{
                \texttt{\tiny{[GM$\rangle$P2]}}
                \texttt{Bogadh 'S' go (4,3) go rathúil. Seo mar atá do ghreille nuashonraithe:} \\
\texttt{```} \\
\texttt{ 1234567} \\
\texttt{╔══╤══╤═╗} \\
\texttt{║◌◌│◌◌│T║ 1} \\
\texttt{╟──┼──┴─╢ 2} \\
\texttt{║◌◌│S◌◌P║ 3} \\
\texttt{╟──┘◌◌MO║ 4} \\
\texttt{║D◌◌◌◌◌◌║ 5} \\
\texttt{║◌◌┌─┬──╢ 6} \\
\texttt{║◌◌│A│◌◌║ 7} \\
\texttt{╚══╧═╧══╝} \\
\texttt{```} \\
\texttt{Tá sibh ag imirt babhta 9 faoi láthair as uasmhéid 28.} \\
\texttt{Tá 15 as 16 pionós carntha agaibh le chéile.} \\
\texttt{Bhog an t{-}imreoir eile réad ar a ngreille.} \\
\\ 
\\ 
\texttt{Cad é do chéad ordú eile?} \\
            }
        }
    }
    & \\ \\

    \theutterance \stepcounter{utterance}  
    & & & \multicolumn{4}{p{0.6\linewidth}}{
        \cellcolor[rgb]{0.9,0.9,0.9}{
            \makecell[{{p{\linewidth}}}]{
                \texttt{\tiny{[GM$\langle$P2]}}
                \texttt{ABAIR: Táimid gar don chríoch. Tá S ailínithe ag (4,3) agus D ag (1,5) ar an dá thaobh. Anois, caithfimid P, M, A, agus T a ailíniú.} \\
\texttt{Mo staid reatha: T(7,1), S(4,3), P(5,3), M(6,4), D(1,5), A(3,7).} \\
\texttt{Do staid reatha (de réir an chomhroinn roimhe): P(1,1), M(1,3), S(2,3{-}>4,3?), O(3,3), A(3,4), D(1,5), T(1,7).} \\
\texttt{Nóta: Dúirt tú go mbog tú O go (3,3). Dúirt tú go raibh S ag (2,3) agat, ach anois tá S ag (4,3) agam. An bhfuil S ag (4,3) agat freisin? Má tá, tá S ailínithe.} \\
\texttt{Molaim an chéad chéim eile:} \\
\texttt{1. Bog tú P ó (1,1) go (1,2) nó (2,1) chun bealach a ghlanadh? Ní hea, ba chóir dúinn P a bhogadh go dtí suíomh atá comhionann liom. Tá P agam ag (5,3). Tá sé sin i bhfad.} \\
\texttt{Ina ionad sin, molaim go n{-}úsáidfimid comhordanáidí níos simplí. Molaim go mbogaimid gach rud go dtí an chéad chill folamh atá in aice leis an mbunús (1,1) nó (1,2) etc. Ach ar dtús, caithfimid a fháil amach cén suíomh atá folamh go deimhin.} \\
\texttt{Mo staid: T(7,1), S(4,3), P(5,3), M(6,4), D(1,5), A(3,7).} \\
\texttt{Do staid: P(1,1), M(1,3), S(4,3), O(3,3), A(3,4), D(1,5), T(1,7).} \\
\texttt{Tá D ailínithe ag (1,5).} \\
\texttt{Tá S ailínithe ag (4,3) (má bhog tú S ansin freisin).} \\
\texttt{Tá O agat ag (3,3). Níl O agam. Caithfidh mé O a fháil? Ní hea, níl O agam ar mo ghreille. Tá sé sin mícheart. Tá O agam ar mo ghreille? Seiceáilim mo ghreille: Níl 'O' le feiceáil. Tá 'M' ag (6,4). Tá 'P' ag (5,3). Tá 'S' ag (4,3). Tá 'T' ag (7,1). Tá 'D' ag (1,5). Tá 'A' ag (3,7). Níl 'O' agam.} \\
\texttt{Ach dúirt tú go raibh 'O' agat ag (4,3) roimhe seo, agus ansin go mbog tú go (3,3).} \\
\texttt{Tá sé seo an{-}chasta. Molaim go n{-}úsáidfimid comhordanáidí nua.} \\
\texttt{Molaim:} \\
\texttt{1. Bog tú P ó (1,1) go (1,2).} \\
\texttt{2. Bog tú M ó (1,3) go (1,1).} \\
\texttt{3. Bog tú A ó (3,4) go (3,3) (ach tá O ansin).} \\
\texttt{Ní hea, tá O ag (3,3).} \\
\texttt{Molaim:} \\
\texttt{1. Bog tú O ó (3,3) go (2,3).} \\
\texttt{2. Ansin, bog tú S ó (4,3) go (3,3).} \\
\texttt{3. Ansin, bog tú P ó (1,1) go (1,2).} \\
\texttt{4. Ansin, bog tú M ó (1,3) go (1,1).} \\
\texttt{5. Ansin, bog tú A ó (3,4) go (4,3) (ach tá S ansin).} \\
\texttt{Tá sé seo ró{-}chasta.} \\
\texttt{Molaim go n{-}úsáidfimid comhordanáidí níos simplí.} \\
\texttt{Molaim go mbogaimid gach rud go dtí an chéad chill folamh atá in aice leis an mbunús (1,1) nó (1,2) etc.} \\
\texttt{Ach ar dtús, caithfimid a fháil amach cén suíomh atá folamh go deimhin.} \\
\texttt{Mo staid: T(7,1), S(4,3), P(5,3), M(6,4), D(1,5), A(3,7).} \\
\texttt{Do staid: P(1,1), M(1,3), S(4,3), O(3,3), A(3,4), D(1,5), T(1,7).} \\
\texttt{Tá D ailínithe ag (1,5).} \\
\texttt{Tá S ailínithe ag (4,3) (má bhog tú S ansin freisin).} \\
\texttt{Tá O agat ag (3,3). Níl O agam. Caithfidh mé O a fháil? Ní hea, níl O agam ar mo ghreille. Tá sé sin mícheart. Tá O agam ar mo ghreille? Seiceáilim mo ghreille: Níl 'O' le feiceáil. Tá 'M' ag (6,4). Tá 'P' ag (5,3). Tá 'S' ag (4,3). Tá 'T' ag (7,1). Tá 'D' ag (1,5). Tá 'A' ag (3,7). Níl 'O' agam.} \\
\texttt{Ach dúirt tú go raibh 'O' agat ag (4,3) roimhe seo, agus ansin go mbog tú go (3,3).} \\
\texttt{Tá sé seo an{-}chasta. Molaim go n{-}úsáidfimid comhordanáidí nua.} \\
\texttt{Molaim:} \\
\texttt{1. Bog tú P ó (1,1) go (1,2).} \\
\texttt{2. Bog tú M ó (1,3) go (1,1).} \\
\texttt{3. Bog tú A ó (3,4) go (3,3) (ach tá O ansin).} \\
\texttt{Ní hea, tá O ag (3,3).} \\
\texttt{Molaim:} \\
\texttt{1. Bog tú O ó (3,3) go (2,3).} \\
\texttt{2. Ansin, bog tú S ó (4,3) go (3,3).} \\
\texttt{3. Ansin, bog tú P ó (1,1) go (1,2).} \\
\texttt{4. Ansin, bog tú M ó (1,3) go (1,1).} \\
\texttt{5. Ansin, bog tú A ó (3,4) go (4,3) (ach tá S ansin).} \\
\texttt{Tá sé seo ró{-}chasta.} \\
\texttt{Molaim go n{-}úsáidfimid comhordanáidí níos simplí.} \\
\texttt{Molaim go mbogaimid gach rud go dtí an chéad chill folamh atá in aice leis an mbunús (1,1) nó (1,2) etc.} \\
\texttt{Ach ar dtús, caithfimid a fháil amach cén suíomh atá folamh go deimhin.} \\
\texttt{Mo staid: T(7,1), S(4,3), P(5,3), M(6,4), D(1,5), A(3,7).} \\
\texttt{Do staid: P(1,1), M(1,3), S(4,3), O(3,3), A(3,4), D(1,5), T(1,7).} \\
\texttt{Tá D ailínithe ag (1,5).} \\
\texttt{Tá S ailínithe ag (4,3) (má bhog tú S ansin freisin).} \\
\texttt{Tá O agat ag (3,3). Níl O agam. Caithfidh mé O a fháil? Ní hea, níl O agam ar mo ghreille. Tá sé sin mícheart. Tá O agam ar mo ghreille? Seiceáilim mo ghreille: Níl 'O' le feiceáil. Tá 'M' ag (6,4). Tá 'P' ag (5,3). Tá 'S' ag (4,3). Tá 'T' ag (7,1). Tá 'D' ag (1,5). Tá 'A' ag (3,7). N} \\
            }
        }
    }
     \\ \\

    \theutterance \stepcounter{utterance}  
    & & & \multicolumn{2}{p{0.3\linewidth}}{
        \cellcolor[rgb]{0.9,0.9,0.9}{
            \makecell[{{p{\linewidth}}}]{
                \texttt{\tiny{[GM$|$GM]}}
                \texttt{ABAIR: Táimid gar don chríoch. Tá S ailínithe ag (4,3) agus D ag (1,5) ar an dá thaobh. Anois, caithfimid P, M, A, agus T a ailíniú.} \\
\texttt{Mo staid reatha: T(7,1), S(4,3), P(5,3), M(6,4), D(1,5), A(3,7).} \\
\texttt{Do staid reatha (de réir an chomhroinn roimhe): P(1,1), M(1,3), S(2,3{-}>4,3?), O(3,3), A(3,4), D(1,5), T(1,7).} \\
\texttt{Nóta: Dúirt tú go mbog tú O go (3,3). Dúirt tú go raibh S ag (2,3) agat, ach anois tá S ag (4,3) agam. An bhfuil S ag (4,3) agat freisin? Má tá, tá S ailínithe.} \\
\texttt{Molaim an chéad chéim eile:} \\
\texttt{1. Bog tú P ó (1,1) go (1,2) nó (2,1) chun bealach a ghlanadh? Ní hea, ba chóir dúinn P a bhogadh go dtí suíomh atá comhionann liom. Tá P agam ag (5,3). Tá sé sin i bhfad.} \\
\texttt{Ina ionad sin, molaim go n{-}úsáidfimid comhordanáidí níos simplí. Molaim go mbogaimid gach rud go dtí an chéad chill folamh atá in aice leis an mbunús (1,1) nó (1,2) etc. Ach ar dtús, caithfimid a fháil amach cén suíomh atá folamh go deimhin.} \\
\texttt{Mo staid: T(7,1), S(4,3), P(5,3), M(6,4), D(1,5), A(3,7).} \\
\texttt{Do staid: P(1,1), M(1,3), S(4,3), O(3,3), A(3,4), D(1,5), T(1,7).} \\
\texttt{Tá D ailínithe ag (1,5).} \\
\texttt{Tá S ailínithe ag (4,3) (má bhog tú S ansin freisin).} \\
\texttt{Tá O agat ag (3,3). Níl O agam. Caithfidh mé O a fháil? Ní hea, níl O agam ar mo ghreille. Tá sé sin mícheart. Tá O agam ar mo ghreille? Seiceáilim mo ghreille: Níl 'O' le feiceáil. Tá 'M' ag (6,4). Tá 'P' ag (5,3). Tá 'S' ag (4,3). Tá 'T' ag (7,1). Tá 'D' ag (1,5). Tá 'A' ag (3,7). Níl 'O' agam.} \\
\texttt{Ach dúirt tú go raibh 'O' agat ag (4,3) roimhe seo, agus ansin go mbog tú go (3,3).} \\
\texttt{Tá sé seo an{-}chasta. Molaim go n{-}úsáidfimid comhordanáidí nua.} \\
\texttt{Molaim:} \\
\texttt{1. Bog tú P ó (1,1) go (1,2).} \\
\texttt{2. Bog tú M ó (1,3) go (1,1).} \\
\texttt{3. Bog tú A ó (3,4) go (3,3) (ach tá O ansin).} \\
\texttt{Ní hea, tá O ag (3,3).} \\
\texttt{Molaim:} \\
\texttt{1. Bog tú O ó (3,3) go (2,3).} \\
\texttt{2. Ansin, bog tú S ó (4,3) go (3,3).} \\
\texttt{3. Ansin, bog tú P ó (1,1) go (1,2).} \\
\texttt{4. Ansin, bog tú M ó (1,3) go (1,1).} \\
\texttt{5. Ansin, bog tú A ó (3,4) go (4,3) (ach tá S ansin).} \\
\texttt{Tá sé seo ró{-}chasta.} \\
\texttt{Molaim go n{-}úsáidfimid comhordanáidí níos simplí.} \\
\texttt{Molaim go mbogaimid gach rud go dtí an chéad chill folamh atá in aice leis an mbunús (1,1) nó (1,2) etc.} \\
\texttt{Ach ar dtús, caithfimid a fháil amach cén suíomh atá folamh go deimhin.} \\
\texttt{Mo staid: T(7,1), S(4,3), P(5,3), M(6,4), D(1,5), A(3,7).} \\
\texttt{Do staid: P(1,1), M(1,3), S(4,3), O(3,3), A(3,4), D(1,5), T(1,7).} \\
\texttt{Tá D ailínithe ag (1,5).} \\
\texttt{Tá S ailínithe ag (4,3) (má bhog tú S ansin freisin).} \\
\texttt{Tá O agat ag (3,3). Níl O agam. Caithfidh mé O a fháil? Ní hea, níl O agam ar mo ghreille. Tá sé sin mícheart. Tá O agam ar mo ghreille? Seiceáilim mo ghreille: Níl 'O' le feiceáil. Tá 'M' ag (6,4). Tá 'P' ag (5,3). Tá 'S' ag (4,3). Tá 'T' ag (7,1). Tá 'D' ag (1,5). Tá 'A' ag (3,7). Níl 'O' agam.} \\
\texttt{Ach dúirt tú go raibh 'O' agat ag (4,3) roimhe seo, agus ansin go mbog tú go (3,3).} \\
\texttt{Tá sé seo an{-}chasta. Molaim go n{-}úsáidfimid comhordanáidí nua.} \\
\texttt{Molaim:} \\
\texttt{1. Bog tú P ó (1,1) go (1,2).} \\
\texttt{2. Bog tú M ó (1,3) go (1,1).} \\
\texttt{3. Bog tú A ó (3,4) go (3,3) (ach tá O ansin).} \\
\texttt{Ní hea, tá O ag (3,3).} \\
\texttt{Molaim:} \\
\texttt{1. Bog tú O ó (3,3) go (2,3).} \\
\texttt{2. Ansin, bog tú S ó (4,3) go (3,3).} \\
\texttt{3. Ansin, bog tú P ó (1,1) go (1,2).} \\
\texttt{4. Ansin, bog tú M ó (1,3) go (1,1).} \\
\texttt{5. Ansin, bog tú A ó (3,4) go (4,3) (ach tá S ansin).} \\
\texttt{Tá sé seo ró{-}chasta.} \\
\texttt{Molaim go n{-}úsáidfimid comhordanáidí níos simplí.} \\
\texttt{Molaim go mbogaimid gach rud go dtí an chéad chill folamh atá in aice leis an mbunús (1,1) nó (1,2) etc.} \\
\texttt{Ach ar dtús, caithfimid a fháil amach cén suíomh atá folamh go deimhin.} \\
\texttt{Mo staid: T(7,1), S(4,3), P(5,3), M(6,4), D(1,5), A(3,7).} \\
\texttt{Do staid: P(1,1), M(1,3), S(4,3), O(3,3), A(3,4), D(1,5), T(1,7).} \\
\texttt{Tá D ailínithe ag (1,5).} \\
\texttt{Tá S ailínithe ag (4,3) (má bhog tú S ansin freisin).} \\
\texttt{Tá O agat ag (3,3). Níl O agam. Caithfidh mé O a fháil? Ní hea, níl O agam ar mo ghreille. Tá sé sin mícheart. Tá O agam ar mo ghreille? Seiceáilim mo ghreille: Níl 'O' le feiceáil. Tá 'M' ag (6,4). Tá 'P' ag (5,3). Tá 'S' ag (4,3). Tá 'T' ag (7,1). Tá 'D' ag (1,5). Tá 'A' ag (3,7). N} \\
            }
        }
    }
    & & \\ \\

    \theutterance \stepcounter{utterance}  
    & & & \multicolumn{2}{p{0.3\linewidth}}{
        \cellcolor[rgb]{0.9,0.9,0.9}{
            \makecell[{{p{\linewidth}}}]{
                \texttt{\tiny{[GM$|$GM]}}
                \texttt{Invalid format: tail is not empty: '} \\
\texttt{Mo staid reatha: T(7,1), S(4,3), P(5,3), M(6,4), D(1,5), A(3,7).} \\
\texttt{Do staid reatha (de réir an chomhroinn roimhe): P(1,1), M(1,3), S(2,3{-}>4,3?), O(3,3), A(3,4), D(1,5), T(1,7).} \\
\texttt{Nóta: Dúirt tú go mbog tú O go (3,3). Dúirt tú go raibh S ag (2,3) agat, ach anois tá S ag (4,3) agam. An bhfuil S ag (4,3) agat freisin? Má tá, tá S ailínithe.} \\
\texttt{Molaim an chéad chéim eile:} \\
\texttt{1. Bog tú P ó (1,1) go (1,2) nó (2,1) chun bealach a ghlanadh? Ní hea, ba chóir dúinn P a bhogadh go dtí suíomh atá comhionann liom. Tá P agam ag (5,3). Tá sé sin i bhfad.} \\
\texttt{Ina ionad sin, molaim go n{-}úsáidfimid comhordanáidí níos simplí. Molaim go mbogaimid gach rud go dtí an chéad chill folamh atá in aice leis an mbunús (1,1) nó (1,2) etc. Ach ar dtús, caithfimid a fháil amach cén suíomh atá folamh go deimhin.} \\
\texttt{Mo staid: T(7,1), S(4,3), P(5,3), M(6,4), D(1,5), A(3,7).} \\
\texttt{Do staid: P(1,1), M(1,3), S(4,3), O(3,3), A(3,4), D(1,5), T(1,7).} \\
\texttt{Tá D ailínithe ag (1,5).} \\
\texttt{Tá S ailínithe ag (4,3) (má bhog tú S ansin freisin).} \\
\texttt{Tá O agat ag (3,3). Níl O agam. Caithfidh mé O a fháil? Ní hea, níl O agam ar mo ghreille. Tá sé sin mícheart. Tá O agam ar mo ghreille? Seiceáilim mo ghreille: Níl 'O' le feiceáil. Tá 'M' ag (6,4). Tá 'P' ag (5,3). Tá 'S' ag (4,3). Tá 'T' ag (7,1). Tá 'D' ag (1,5). Tá 'A' ag (3,7). Níl 'O' agam.} \\
\texttt{Ach dúirt tú go raibh 'O' agat ag (4,3) roimhe seo, agus ansin go mbog tú go (3,3).} \\
\texttt{Tá sé seo an{-}chasta. Molaim go n{-}úsáidfimid comhordanáidí nua.} \\
\texttt{Molaim:} \\
\texttt{1. Bog tú P ó (1,1) go (1,2).} \\
\texttt{2. Bog tú M ó (1,3) go (1,1).} \\
\texttt{3. Bog tú A ó (3,4) go (3,3) (ach tá O ansin).} \\
\texttt{Ní hea, tá O ag (3,3).} \\
\texttt{Molaim:} \\
\texttt{1. Bog tú O ó (3,3) go (2,3).} \\
\texttt{2. Ansin, bog tú S ó (4,3) go (3,3).} \\
\texttt{3. Ansin, bog tú P ó (1,1) go (1,2).} \\
\texttt{4. Ansin, bog tú M ó (1,3) go (1,1).} \\
\texttt{5. Ansin, bog tú A ó (3,4) go (4,3) (ach tá S ansin).} \\
\texttt{Tá sé seo ró{-}chasta.} \\
\texttt{Molaim go n{-}úsáidfimid comhordanáidí níos simplí.} \\
\texttt{Molaim go mbogaimid gach rud go dtí an chéad chill folamh atá in aice leis an mbunús (1,1) nó (1,2) etc.} \\
\texttt{Ach ar dtús, caithfimid a fháil amach cén suíomh atá folamh go deimhin.} \\
\texttt{Mo staid: T(7,1), S(4,3), P(5,3), M(6,4), D(1,5), A(3,7).} \\
\texttt{Do staid: P(1,1), M(1,3), S(4,3), O(3,3), A(3,4), D(1,5), T(1,7).} \\
\texttt{Tá D ailínithe ag (1,5).} \\
\texttt{Tá S ailínithe ag (4,3) (má bhog tú S ansin freisin).} \\
\texttt{Tá O agat ag (3,3). Níl O agam. Caithfidh mé O a fháil? Ní hea, níl O agam ar mo ghreille. Tá sé sin mícheart. Tá O agam ar mo ghreille? Seiceáilim mo ghreille: Níl 'O' le feiceáil. Tá 'M' ag (6,4). Tá 'P' ag (5,3). Tá 'S' ag (4,3). Tá 'T' ag (7,1). Tá 'D' ag (1,5). Tá 'A' ag (3,7). Níl 'O' agam.} \\
\texttt{Ach dúirt tú go raibh 'O' agat ag (4,3) roimhe seo, agus ansin go mbog tú go (3,3).} \\
\texttt{Tá sé seo an{-}chasta. Molaim go n{-}úsáidfimid comhordanáidí nua.} \\
\texttt{Molaim:} \\
\texttt{1. Bog tú P ó (1,1) go (1,2).} \\
\texttt{2. Bog tú M ó (1,3) go (1,1).} \\
\texttt{3. Bog tú A ó (3,4) go (3,3) (ach tá O ansin).} \\
\texttt{Ní hea, tá O ag (3,3).} \\
\texttt{Molaim:} \\
\texttt{1. Bog tú O ó (3,3) go (2,3).} \\
\texttt{2. Ansin, bog tú S ó (4,3) go (3,3).} \\
\texttt{3. Ansin, bog tú P ó (1,1) go (1,2).} \\
\texttt{4. Ansin, bog tú M ó (1,3) go (1,1).} \\
\texttt{5. Ansin, bog tú A ó (3,4) go (4,3) (ach tá S ansin).} \\
\texttt{Tá sé seo ró{-}chasta.} \\
\texttt{Molaim go n{-}úsáidfimid comhordanáidí níos simplí.} \\
\texttt{Molaim go mbogaimid gach rud go dtí an chéad chill folamh atá in aice leis an mbunús (1,1) nó (1,2) etc.} \\
\texttt{Ach ar dtús, caithfimid a fháil amach cén suíomh atá folamh go deimhin.} \\
\texttt{Mo staid: T(7,1), S(4,3), P(5,3), M(6,4), D(1,5), A(3,7).} \\
\texttt{Do staid: P(1,1), M(1,3), S(4,3), O(3,3), A(3,4), D(1,5), T(1,7).} \\
\texttt{Tá D ailínithe ag (1,5).} \\
\texttt{Tá S ailínithe ag (4,3) (má bhog tú S ansin freisin).} \\
\texttt{Tá O agat ag (3,3). Níl O agam. Caithfidh mé O a fháil? Ní hea, níl O agam ar mo ghreille. Tá sé sin mícheart. Tá O agam ar mo ghreille? Seiceáilim mo ghreille: Níl 'O' le feiceáil. Tá 'M' ag (6,4). Tá 'P' ag (5,3). Tá 'S' ag (4,3). Tá 'T' ag (7,1). Tá 'D' ag (1,5). Tá 'A' ag (3,7). N} \\
            }
        }
    }
    & & \\ \\

    \theutterance \stepcounter{utterance}  
    & & \multicolumn{4}{p{0.6\linewidth}}{
        \cellcolor[rgb]{0.9,0.9,0.9}{
            \makecell[{{p{\linewidth}}}]{
                \texttt{\tiny{[GM$\rangle$P2]}}
                \texttt{Pionós: Ní cheadaítear aon rud a bheith i do theachtaireacht tar éis an ordaithe!} \\
\\ 
\texttt{Cinntigh nach bhfuil i do fhreagra ach `ABAIR: <TEACHTAIREACHT>` nó `BOG: <RUD>, (<X>, <Y>)`, agus tada eile!} \\
\\ 
\texttt{Tá 16 as 16 pionós carntha agaibh le chéile.} \\
\texttt{Déan iarracht arís, le do thoil!} \\
            }
        }
    }
    & \\ \\

    \theutterance \stepcounter{utterance}  
    & & & \multicolumn{4}{p{0.6\linewidth}}{
        \cellcolor[rgb]{0.9,0.9,0.9}{
            \makecell[{{p{\linewidth}}}]{
                \texttt{\tiny{[GM$\langle$P2]}}
                \texttt{ABAIR: Táimid tar éis 16 pionós a chur le chéile. Caillimid an cluiche. Críochnaímid anseo.} \\
            }
        }
    }
     \\ \\

    \theutterance \stepcounter{utterance}  
    & & & \multicolumn{2}{p{0.3\linewidth}}{
        \cellcolor[rgb]{0.9,0.9,0.9}{
            \makecell[{{p{\linewidth}}}]{
                \texttt{\tiny{[GM$|$GM]}}
                \texttt{ABAIR: Táimid tar éis 16 pionós a chur le chéile. Caillimid an cluiche. Críochnaímid anseo.} \\
            }
        }
    }
    & & \\ \\

    \theutterance \stepcounter{utterance}  
    & & \multicolumn{4}{p{0.6\linewidth}}{
        \cellcolor[rgb]{0.9,0.9,0.9}{
            \makecell[{{p{\linewidth}}}]{
                \texttt{\tiny{[P1$\langle$GM]}}
                \texttt{Bogadh 'S' go (4,3) go rathúil. Seo mar atá do ghreille nuashonraithe:} \\
\texttt{```} \\
\texttt{ 1234567} \\
\texttt{╔══╤══╤═╗} \\
\texttt{║◌P│◌◌│◌║ 1} \\
\texttt{╟──┼──┴─╢ 2} \\
\texttt{║M◌│SO◌◌║ 3} \\
\texttt{╟──┘◌A◌◌║ 4} \\
\texttt{║D◌◌◌◌◌◌║ 5} \\
\texttt{║◌◌┌─┬──╢ 6} \\
\texttt{║T◌│◌│◌◌║ 7} \\
\texttt{╚══╧═╧══╝} \\
\texttt{```} \\
\texttt{Tá sibh ag imirt babhta 9 faoi láthair as uasmhéid 28.} \\
\texttt{Tá 16 as 16 pionós carntha agaibh le chéile.} \\
\texttt{Sheol an t{-}imreoir eile an teachtaireacht seo:} \\
\\ 
\\ 
\texttt{"Táimid tar éis 16 pionós a chur le chéile. Caillimid an cluiche. Críochnaímid anseo."} \\
\\ 
\\ 
\texttt{Cad é do chéad ordú eile?} \\
            }
        }
    }
    & \\ \\

    \theutterance \stepcounter{utterance}  
    & \multicolumn{4}{p{0.6\linewidth}}{
        \cellcolor[rgb]{0.9,0.9,0.9}{
            \makecell[{{p{\linewidth}}}]{
                \texttt{\tiny{[P1$\rangle$GM]}}
                \texttt{ABAIR: críochnaithe!} \\
            }
        }
    }
    & & \\ \\

    \theutterance \stepcounter{utterance}  
    & & & \multicolumn{2}{p{0.3\linewidth}}{
        \cellcolor[rgb]{0.9,0.9,0.9}{
            \makecell[{{p{\linewidth}}}]{
                \texttt{\tiny{[GM$|$GM]}}
                \texttt{ABAIR: críochnaithe!} \\
            }
        }
    }
    & & \\ \\

    \theutterance \stepcounter{utterance}  
    & & \multicolumn{4}{p{0.6\linewidth}}{
        \cellcolor[rgb]{0.9,0.9,0.9}{
            \makecell[{{p{\linewidth}}}]{
                \texttt{\tiny{[GM$\rangle$P2]}}
                \texttt{Cuireadh do theachtaireacht ar aghaidh chuig an imreoir eile.} \\
\\ 
\\ 
\texttt{Tá sibh ag imirt babhta 10 faoi láthair as uasmhéid 28.} \\
\texttt{Tá 16 as 16 pionós carntha agaibh le chéile.} \\
\texttt{Sheol an t{-}imreoir eile an teachtaireacht seo:} \\
\\ 
\\ 
\texttt{"críochnaithe!"} \\
\\ 
\\ 
\texttt{Cad é do chéad ordú eile?} \\
            }
        }
    }
    & \\ \\

    \theutterance \stepcounter{utterance}  
    & & & \multicolumn{4}{p{0.6\linewidth}}{
        \cellcolor[rgb]{0.9,0.9,0.9}{
            \makecell[{{p{\linewidth}}}]{
                \texttt{\tiny{[GM$\langle$P2]}}
                \texttt{ABAIR: críochnaithe!} \\
            }
        }
    }
     \\ \\

    \theutterance \stepcounter{utterance}  
    & & & \multicolumn{2}{p{0.3\linewidth}}{
        \cellcolor[rgb]{0.9,0.9,0.9}{
            \makecell[{{p{\linewidth}}}]{
                \texttt{\tiny{[GM$|$GM]}}
                \texttt{ABAIR: críochnaithe!} \\
            }
        }
    }
    & & \\ \\

    \theutterance \stepcounter{utterance}  
    & & & \multicolumn{2}{p{0.3\linewidth}}{
        \cellcolor[rgb]{0.9,0.9,0.9}{
            \makecell[{{p{\linewidth}}}]{
                \texttt{\tiny{[GM$|$GM]}}
                \texttt{true} \\
            }
        }
    }
    & & \\ \\

    \theutterance \stepcounter{utterance}  
    & & & \multicolumn{2}{p{0.3\linewidth}}{
        \cellcolor[rgb]{0.9,0.9,0.9}{
            \makecell[{{p{\linewidth}}}]{
                \texttt{\tiny{[GM$|$GM]}}
                \texttt{Initial states:} \\
\texttt{```} \\
\texttt{Player 1:} \\
\texttt{ 1234567} \\
\texttt{╔══╤══╤═╗} \\
\texttt{║◌P│◌M│◌║ 1} \\
\texttt{╟──┼──┴─╢ 2} \\
\texttt{║◌◌│◌S◌O║ 3} \\
\texttt{╟──┘◌A◌◌║ 4} \\
\texttt{║D◌◌◌◌◌◌║ 5} \\
\texttt{║◌◌┌─┬──╢ 6} \\
\texttt{║T◌│◌│◌◌║ 7} \\
\texttt{╚══╧═╧══╝} \\
\\ 
\\ 
\texttt{Player 2:} \\
\texttt{ 1234567} \\
\texttt{╔══╤══╤═╗} \\
\texttt{║◌◌│◌◌│T║ 1} \\
\texttt{╟──┼──┴─╢ 2} \\
\texttt{║◌◌│◌S◌P║ 3} \\
\texttt{╟──┘◌◌MO║ 4} \\
\texttt{║D◌◌◌◌◌◌║ 5} \\
\texttt{║◌◌┌─┬──╢ 6} \\
\texttt{║◌◌│A│◌◌║ 7} \\
\texttt{╚══╧═╧══╝} \\
\\ 
\texttt{```} \\
            }
        }
    }
    & & \\ \\

    \theutterance \stepcounter{utterance}  
    & & & \multicolumn{2}{p{0.3\linewidth}}{
        \cellcolor[rgb]{0.9,0.9,0.9}{
            \makecell[{{p{\linewidth}}}]{
                \texttt{\tiny{[GM$|$GM]}}
                \texttt{End states:} \\
\texttt{```} \\
\texttt{Player 1:} \\
\texttt{ 1234567} \\
\texttt{╔══╤══╤═╗} \\
\texttt{║◌P│◌◌│◌║ 1} \\
\texttt{╟──┼──┴─╢ 2} \\
\texttt{║M◌│SO◌◌║ 3} \\
\texttt{╟──┘◌A◌◌║ 4} \\
\texttt{║D◌◌◌◌◌◌║ 5} \\
\texttt{║◌◌┌─┬──╢ 6} \\
\texttt{║T◌│◌│◌◌║ 7} \\
\texttt{╚══╧═╧══╝} \\
\\ 
\\ 
\texttt{Player 2:} \\
\texttt{ 1234567} \\
\texttt{╔══╤══╤═╗} \\
\texttt{║◌◌│◌◌│T║ 1} \\
\texttt{╟──┼──┴─╢ 2} \\
\texttt{║◌◌│S◌◌P║ 3} \\
\texttt{╟──┘◌◌MO║ 4} \\
\texttt{║D◌◌◌◌◌◌║ 5} \\
\texttt{║◌◌┌─┬──╢ 6} \\
\texttt{║◌◌│A│◌◌║ 7} \\
\texttt{╚══╧═╧══╝} \\
\\ 
\texttt{```} \\
            }
        }
    }
    & & \\ \\

    \theutterance \stepcounter{utterance}  
    & & & \multicolumn{2}{p{0.3\linewidth}}{
        \cellcolor[rgb]{0.9,0.9,0.9}{
            \makecell[{{p{\linewidth}}}]{
                \texttt{\tiny{[GM$|$GM]}}
                \texttt{* success: True} \\
\texttt{* lose: False} \\
\texttt{* aborted: False} \\
\texttt{{-}{-}{-}{-}{-}{-}{-}} \\
\texttt{* Shifts: 3.00} \\
\texttt{* Max Shifts: 12.00} \\
\texttt{* Min Shifts: 6.00} \\
\texttt{* End Distance Sum: 24.37} \\
\texttt{* Init Distance Sum: 21.20} \\
\texttt{* Expected Distance Sum: 29.33} \\
\texttt{* Penalties: 16.00} \\
\texttt{* Max Penalties: 16.00} \\
\texttt{* Rounds: 10.00} \\
\texttt{* Max Rounds: 28.00} \\
\texttt{* Object Count: 7.00} \\
            }
        }
    }
    & & \\ \\

\end{supertabular}
}

\end{document}
